\documentclass[11pt]{article}
\usepackage[T1]{fontenc}
\usepackage[utf8]{inputenc}
\usepackage{lmodern}
\usepackage[margin=1in]{geometry}
\usepackage{amsmath,amssymb,amsthm,mathtools,bm}
\usepackage{booktabs,array,longtable}
\usepackage{enumitem}
\usepackage{hyperref}
\usepackage{xurl}
\usepackage{seqsplit}
\newcommand{\OPHPaperReleaseID}{r2055}
\newcommand{\OPHPaperReleaseDate}{October 28, 2026}
\newcommand{\OPHPaperReleaseBanner}{%
\begin{center}
\small\textbf{Paper release:} \texttt{\OPHPaperReleaseID}\quad\textbf{Released:} \OPHPaperReleaseDate
\end{center}
}
\newcommand{\PaperReleaseID}{\OPHPaperReleaseID}
\newcommand{\PaperReleaseDate}{\OPHPaperReleaseDate}
\newcommand{\PaperReleaseBanner}{%
\begin{center}
\small\textbf{Paper release:} \texttt{\PaperReleaseID}\quad\textbf{Released:} \PaperReleaseDate
\end{center}
}


\hypersetup{colorlinks=true,linkcolor=blue,citecolor=blue,urlcolor=blue}
\hypersetup{pdftitle={The Particle Zoo under Declared Matter Premises: One Generation, Interaction Channels and Mass Diagnostics},pdfauthor={Bernhard Mueller; Mario Poneder; Alexander Osika; Kai Xue; Daniel Ari Friedman}}
\setlength{\emergencystretch}{3em}
\setlength{\LTleft}{0pt}
\setlength{\LTright}{0pt}
\setlength{\tabcolsep}{5pt}
\renewcommand{\arraystretch}{1.08}

\newtheorem{theorem}{Theorem}[section]
\newtheorem{corollary}[theorem]{Corollary}
\newtheorem{proposition}[theorem]{Proposition}
\newtheorem{lemma}[theorem]{Lemma}
\newtheorem{definition}[theorem]{Definition}
\newtheorem{axiom}[theorem]{Axiom}
\newtheorem{assumption}[theorem]{Assumption}
\newtheorem{remark}[theorem]{Remark}

\newcommand{\SU}{\mathrm{SU}}
\newcommand{\U}{\mathrm{U}}
\newcommand{\ellbar}{\bar{\ell}}
\newcommand{\Conf}{\mathrm{Conf}}
\newcommand{\ophid}[1]{\path{#1}}
\newcommand{\decnum}[1]{\vcenter{\hbox{\parbox{0.68\linewidth}{\small\ttfamily\seqsplit{#1}}}}}

\title{The Particle Zoo under Declared Matter Premises:\\
One Generation, Interaction Channels and Mass Diagnostics}
\author{B.\ M\"uller \and Mario Poneder \and Alexander Osika \and Kai Xue \and Daniel Ari Friedman}
\date{\OPHPaperReleaseDate}

\usepackage{microtype}
\begin{document}
\maketitle
\OPHPaperReleaseBanner

\begin{center}
\small\textbf{Author affiliations:} Bernhard Mueller, Pragma Research Inc.; Alexander Osika, EtherWorks; Daniel Ari Friedman, Active Inference Institute and FractiAI.
\end{center}

\begin{abstract}
A particle spectrum requires more than a gauge representation: it requires
physical fields, their dynamics, a vacuum and a mass observable. We separate
these requirements for finite Observer Patch Holography constructions. A
specified twelve-current response, with a declared spin lift and exterior
module, gives the fifteen internal states of one Standard Model generation
and its three charged-fermion Higgs channels. Their family coefficients are
free. Calibrated finite actions permit observer-record recovery of supplied
mass parameters; the tested uncalibrated repair prescriptions fail the charged
hierarchies. We give exact normalization bounds and restricted
nonidentifiability results explaining those failures. Conditional flavor,
electroweak, clock and pole constructions identify the additional hypotheses
needed for physical masses. We also distinguish finite record continuation
from particle propagation and state the continuum, kinetic-normalization and
source-selection requirements. The construction establishes conditional
representation and readout results, rather than a source-derived measured
particle spectrum or a complete interacting Standard Model.
\end{abstract}

% Prose definitions shared by the current interpretive and physical papers.
% This panel changes no axiom, theorem hypothesis, or formal claim class.
\paragraph{Interpretation and state variables.}
OPH proposes one structure with two descriptions: a timeless, self-referential
complete history and its intrinsic computation through local events, records,
and responses to admissible interventions. The correspondence must preserve
those dependencies; no external host or additional time in which the history
is played is required. A candidate must specify the history domain and closure
operator; the schematic equation does not construct or select a unique universe.
An actual quantum record history also need not
assign simultaneous outcomes to incompatible unperformed measurements.

The proposed primitives are patch states, records, interfaces, and admissible
operations. Physical space, duration, and matter are intended emergent readouts.
A supplied carrier incidence, geometric domain, action, representation, or
clock remains a premise of the construction using it. Complete-history
closure, an instantaneous repair fixed point, a public normal form, and a
stationary probability law are different objects. A stationary law can
support changing individual histories. An operational observer is defined by
its access and record interfaces; consciousness requires an additional
psychophysical interpretation.

\paragraph{Time conventions.}
Host runtime, solver counters, event precedence or logical numbering, model
coordinates or lapse, modular parameters, and calibrated observer-clock
readings have distinct roles. Their identification requires explicit maps.
A modular parameter is dimensionless; its geometric and operational readings
require the stated branch and clock premises. Consistent local clock
comparisons require no preferred universal clock or identical clock rates.
Shared time parameters in restricted benchmarks remain declared model choices.
Experienced succession and a final instantaneous cosmic endpoint do not follow
from a modular-flow or finite-convergence theorem.


\section{Particle structure and physical mass}

Observer Patch Holography (OPH) models observations through finite accessible
algebras, retained records and local repair. The regulated observer
construction provides explicit source operators and instruments
\cite{ophmodel,ophrepairlaw}. This paper asks which particle statements follow
from those operators and which require an additional field-theory attachment.
The gauge classification is developed in Ref.~\cite{ophGauge}; the finite
matter decoder is restated below so that its hypotheses and its limitations
can be assessed without an inferred identification of internal states with
laboratory particles.

The principal conclusions are:
\begin{enumerate}[leftmargin=2em]
\item A specified closed twelve-current law admits a conditional
Standard Model matter and Higgs representation. Its symmetry identities
permit arbitrary family Yukawa matrices.
\item Calibrated actions realize charged-matter structure and permit recovery
of their supplied mass parameters from simulated observer records. This is a
realization and readout test, not a derivation of the calibration inputs.
\item The uncalibrated response laws tested here do not reproduce the measured
charged spectrum. Exact exclusions apply to stated prescriptions, including
normalized projector averages and an affine family of quadratic responses.
\item The physical chiral interacting continuum, source-selected vacuum,
kinetic forms, Yukawa coefficients and physical mass/clock comparison are
additional requirements. No measured particle spectrum is derived.
\end{enumerate}

A finite internal eigenvalue may be a representation label, a memory level,
a relaxation rate, a Hamiltonian coefficient or a mass. Those meanings require
different observation maps. In particular, a repair iteration counter is not
an operational clock, and a finite relaxation gap is not a relativistic
mass-shell pole. The word \emph{source-derived} below requires the complete
source program to determine an object independently of its empirical target,
including the ancestry of its discrete choices and coefficients. Absence of
a target file at runtime does not establish this independence.

A particle identification must begin with one excitation or sector of one
continuing source, and preserve that object through its charge, propagation,
detector, kinetic and interaction maps. The different representation labels,
memory levels, decay rates and benchmark field poles below cannot be pooled
into a spectrum by choosing the closest numerical row. For a stable excitation
in an appropriate vacuum regime, nonzero response residue and dispersion
must accompany the quantum numbers; unstable, confined and infrared charged
sectors require their own observable definitions. A topological-defect
interpretation additionally needs an actual invariant of the admitted native
updates. Nonzero repair frustration or resource expenditure alone does not
identify inertial mass. The source must connect the same excitation's excess
energy and clocked propagation, including kinetic normalization. The current
construction contains no established native defect-to-physical-mass map.

\section{Source and attachment premises}\label{sec:particle-input-ledger}

% Shared panel of the theorem premises A1--A3. Included by papers that state
% the premises in full. Canonical machine source: claims/axiom_registry.yaml;
% reviewed reference: docs/AXIOM_REFERENCE.md. Do not hand-copy a competing
% formal block into a paper; include this fragment.

\paragraph{A1 (oriented observer-patch federation and spherical support).}
There exists an observer patch net on an oriented spherical screen: at every
finite resolution each local carrier has twelve primitive boundary ports
forming the vertices of an oriented triangular boundary with 30 edges and 20
faces, combinatorially the boundary of an icosahedron, and carriers join
through typed seams and coherent triple overlaps, refine to an oriented
spherical support, and expose local state, readback, records, repair moves,
and checkpoints. Their complete infinitesimal port response is represented
faithfully on a finite-dimensional unitary response space and is closed under
the commutators generated by ordered response composition. Formally, for
every regulator \(r\) in a
directed system there is a typed object
\(\mathfrak N_r=(\mathcal P_r,\mathcal A_r,\mathcal R_r,\mathcal I_r,
\mathcal U_r,\mathcal C_r,N_r,S_r,b_r)\):
a finite patch/overlap category with an isotone local algebra net and
central record algebras; a federation of finite carriers, each with twelve
primitive pairwise-orthogonal central port projections summing to one and a
boundary packet \(K=(P,E,F,o)\) with \(|P|=12\), \(|E|=30\), \(|F|=20\),
degree-five ports, five-cycle links, and coherently oriented two-face edges,
isomorphic to the icosahedral boundary complex with no preferred labels;
seam algebras with unital restrictions and coherent triple-overlap cocycles
forming the nerve \(N_r\); a finite edge-midpoint refinement support
\(S_r\) realized orientation-preservingly in \(S^2\) with mesh tending to
zero; and a source-bound degree-one bridge \(b_r:N_r\to S_r\), all commuting
with refinement. Each carrier also has
\(V_{r,i}=\mathbb R[P_{r,i}]\), a finite-dimensional complex Hilbert space
\(H_{r,i}\), and an injective real-linear map
\[
D_{r,i}:V_{r,i}\longrightarrow\mathfrak u(H_{r,i}).
\]
The image \(\mathfrak g_{r,i}=D_{r,i}(V_{r,i})\) is closed under
commutators and complete for the declared quotient-visible infinitesimal
port response. The primitive port probes span \(V_{r,i}\), no additional
public response direction is omitted, and
\(-\operatorname{Tr}(D_{r,i}(v)D_{r,i}(w))\) is positive definite. The
response construction is natural under admissible presentation equivalence
and refinement. A1 does not select a particular response table,
inverse-port law, Lie type, global group, particle interpretation, or metric
content.

\paragraph{A2 (observer agreement).}
Observers operating on the screen agree on the meaning of the data they
jointly interpret. Formally, the interpretation map
\(\mathcal J_r:\mathsf{Data}_r\to\mathsf{Meaning}_r\) from
observer-accessible data to operational meanings is natural with respect to
every visible overlap restriction, recharting, seam translation,
higher-overlap map, federation map, and refinement map: every declared
data-access diagram for accepted public data commutes after interpretation,
and \(\mathcal J_r\circ c_{s\to r}=C_{s\to r}\circ\mathcal J_s\) across
resolutions. A2 constrains accepted shared data only, and the domain
it quantifies over is supplied by the A1 interfaces rather than by A2
itself.

Accepted reversible overlap transports form a groupoid \(\mathcal O_r\).
For a complete carrier chart \(o\), its closed paths
\(\operatorname{Hol}_r(o)\) map onto the orientation-preserving incidence
automorphisms \(\operatorname{Aut}^{+}(K_{r,i})\). For every proper carrier
automorphism \(a\), A2 supplies a closed path \(\gamma_a\) and one
projective implementer
\([U_a]\in\operatorname{PU}(H_{r,i})\) satisfying
\[
\operatorname{Ad}_{[U_a]}D_{r,i}(v)=D_{r,i}(a\cdot v).
\]
Writing
\[
G_{D,r,i}^{0}
=
\left\langle
\exp(tD_{r,i}(v)):t\in\mathbb R,\ v\in V_{r,i}
\right\rangle^{0},
\]
the implementer is endogenous:
\[
[U_a]=[g_a c_a],
\qquad
g_a\in G_{D,r,i}^{0},
\quad
c_a\in C_{U(H_{r,i})}(\mathfrak g_{r,i}).
\]
Closed paths compose projectively and the construction commutes with
refinement. A unitary lift need not be unique. Scalar phases and centralizer
actions are permitted because they act trivially on the response algebra.
An independent spectator carrying a matching projective action cannot
replace the response-generated factor or contribute a public response
direction. The closed path, port action, and implementer must arise from the
same accepted source history in an executable realization.

The two response clauses support an abstract target-free compact Lie-type
theorem. They do not certify that an executable producer emitted a particular
current fixture. Such a certification must reconstruct the port generators,
commutator, and projective implementers from ordered source histories without
a named current model upstream. A2 does not imply global state
extension, repair termination, confluence, unique normal forms, record
durability, Byzantine safety, dissemination, or commutation of state
optimization with refinement; raw mismatch and repair precede acceptance.

\paragraph{A3 (conditional maximum randomness).}
Everything that observer agreement leaves unconstrained is maximally random.
Formally, at finite regulator \(r\) a state is a compatible family of local
normalized states on the accessible algebra net; \(\mathcal K_r\) is the
nonempty convex set of such families satisfying the finite observer-visible
constraints supplied by A1 and A2 through an A1-generated constraint
grammar with a factorization theorem; and, relative to an exact reference
family \(\tau_r\), an A1-generated observer cover \(\mathcal G_r\), and
strictly positive exact weights \(w_{r,P}\) constructed from quotient-visible
A1 data by a rule natural under admissible presentation equivalence, the
realized state is the
information projection. The cover restriction map is injective on
\(\mathcal K_r\); equivalently, at smooth points no nonzero feasible
tangent is invisible to the whole cover. The realized state is
\[
\rho_r=\operatorname*{arg\,min}_{\rho\in\mathcal K_r}
\sum_{P\in\mathcal G_r} w_{r,P}\,D(\rho_{r,P}\Vert\tau_{r,P}).
\]
The same projection rule applies separately to each declared optimizer
object type (ontic state, inference state, or transition distribution), and
statements about one type do not transfer to another.
This is equivalent to weighted local entropy maximization when every local
reference density is identity-proportional in its declared trace. "Random"
means least informative relative to the
declared reference and the complete agreement constraint set. A3
selects one state inside one fixed feasible space. It does not compare
unrelated state spaces or field contents, does not imply optimizer
compatibility across refinement, and does not imply recovery, alignment, or
mixing; each such statement carries its own theorem, interface, or
countermodel.

The complete formal statement of A1--A3, including its constraint map and
non-implications, is given in the
\href{https://github.com/FloatingPragma/observer-patch-holography/blob/r2053/docs/AXIOM_REFERENCE.md}
{premise reference}.

% Shared closure panel. Included by papers that state the premises A1--A3,
% immediately after PREMISE_BASIS.tex. It states the closure hypothesis
% of docs/BASIS.md, which no theorem takes as a premise. Reviewed
% references: docs/BASIS.md and docs/AXIOM_REFERENCE.md.

\paragraph{The closure hypothesis.}
The theorem premises A1--A3 state the library's hypotheses on the carrier and
its response, on internal transport and on local state selection. The closure hypothesis
states that the Universe is a timeless, self-referential structure whose
observers recover the specification of their own world from their records
and realize it, so that the description and the described system are one
object. Every quantity that
has both a construction-side reading and a readback-side reading must take
the same value once a typed bridge proves that the two readings denote one
invariant. The equality is forced by self-identity. Constructing that bridge
and the return map is part of the physics. Existence, uniqueness, and
stability are separate determinacy tests on the resulting equation. The name
given to such an equation carries no mathematical weight: after the typed
identification, unequal readings would describe two systems rather than the
single self-referential universe.

The closed object is the complete inhabited history, whose internal events
and responses to interventions are its proposed intrinsic computation.
This reading requires no external computer or playback time. It is distinct
from fixedness of a repair state, a public record, or a probability law, and
does not itself construct the history or prove a final cosmic endpoint.

The local screen-grain proposal
\[
P=\varphi+\sqrt{\pi}/A_T(P),
\]
with \(A_T(P)\) the inverse coupling returned by the declared source map,
has one interval-certified fixed point per declared map on its analytic domain.
Interpreting that root as the laboratory coupling additionally requires a
target-independent choice of the physical map, a typed bridge identifying
the two sides as readings of one quantity, and same-scheme hadronic spectral
transport to the Thomson endpoint.
The displayed map's certified fixed point sits at $P=1.6309720959$ and
returns $A_T=136.9948352$; the certified gauge-width variant's fixed
point sits at $P=1.6309682414$ and returns $137.0356601$. Through the
same outer equation the measured 2022 CODATA coupling
$\alpha^{-1}=137.035999177(21)$ reads back the grain value
$P=1.6309682094$. The comparisons stand at $3.0\times10^{-4}$ relative
for the displayed map and at $2.5\times10^{-6}$ for the gauge-width
variant, as diagnostics of the declared source maps. The latter absolute
residual is $3.3904\times10^{-4}$, compared with experimental standard
uncertainty $2.1\times10^{-8}$; this is not agreement at measured precision.
The gauge width supplies 99.2 percent of the gap between the bare root and the
measured value; the residual $3.3904\times10^{-4}$ is addressed to same-scheme
hadronic transport to the Thomson endpoint.

The direct global proposal is
\[
N=\log M_0(\mathfrak U_N).
\]
The bounded completion class defined by base agreement, positivity, and the
carrier bound does not select a unique cosmic value. Universal all-rung
membership in a complete A1--A3 capacity-source contract and the corresponding
executable-to-Lean bridge are absent. The incomplete antecedent determines no
direct \(N\), and no theorem extends the bounded-class conclusion to the
complete source class. A
positive result must complete the source antecedent and prove one physical
zero, either within the A1--A3 source contract or through a separately
named stronger source law. A separate
common-load branch starts from
\[
N_0=\pi\exp\!\left(\frac{6\pi}{P\,\alpha_U(P)}\right).
\]
On the finite collar branch, the declared total reserve expectation \(P/4\)
and six-class equidistribution give presence probability \(P/24\) for each
declared class. If one class is physically selected as the blocked event, its
scalar-weighted presence receipt is discharged, and that normalized
collar-survival factor is proved to act on the global capacity, the
corresponding candidate is
\[
N_{\mathrm{pres}}=N_0\left(1-\frac{P}{24}\right),
\qquad
\ln\frac{N_{\mathrm{pres}}}{\pi}
=
\frac{6\pi}{P\,\alpha_U(P)}
+\ln\left(1-\frac{P}{24}\right).
\]
The exponential alternative \(N_{\mathrm{Pois}}=N_0e^{-P/24}\) requires a
separate mean-count or continuum carrier. Exact neutral, one-class, and
six-class-total actions share the declared local datum, obey positive
composition and finite cut-count regrouping, and disagree globally. The
finite source therefore selects no action or blocked-event semantics. The
source class determines no named-law value, and its horizon branch has no
capacity object. A stronger source-derived action would
require the screen/electroweak bridge, physical common-load bridge, physical
\(\mathbb Z_6\) seam action, scalar-weighted receipt, and horizon-record
identification. On the source-forward numerical
branch,
\[
N_0=3.5321\times10^{122},
\]
so the finite-presence candidate is
\(3.2921\times10^{122}\) and the exponential candidate is
\(3.3001\times10^{122}\). The full weighted Planck
base-\(\Lambda\)CDM chain gives the comparison coordinate
\(3.3129\times10^{122}\). The respective residuals are
\(-0.63\) percent and \(-0.39\) percent. Both comparisons were exposed before
the branch choice and carry no predictive weight.


These premises organize the conditional OPH mathematical framework. They do not select
every quantitative particle completion. The screen-first spherical support,
the local twelve-port carrier and the federation overlap nerve are distinct
objects; identifying their physical carriers consumes the separate incidence,
access, overlap and refinement hypotheses \cite{ophmicrophysics,ophGauge}.

\begin{center}
\small
\begin{tabular}{p{0.23\linewidth}p{0.68\linewidth}}
\toprule
Object & Hypothesis needed beyond its finite algebraic calculation \\
\midrule
Gauge currents & Selection of the complete physically available response,
rather than a chosen control subspace. \\
Matter module & Spin cover, exterior-sector choice, physical spin and CAR
statistics, and a field interpretation. \\
Three families & Attachment of a protected rank-three band to three physical
generations, with specified symmetry breaking. \\
Mass coefficients & A source-selected Higgs response distinguished from the
kinetic metric, a vacuum and a common normalization. \\
Physical spectrum & Controlled chiral continuum, interactions, renormalization,
infrared treatment, observable singularities and an operational clock. \\
\bottomrule
\end{tabular}
\end{center}

The quantitative pixel branch uses a dimensionless closure coordinate \(P\).
An interval contraction certificate establishes a unique root on the declared
numerical domain; it does not identify that root with a physical endpoint.
The distinction is essential: a scalar fixes a mass scale only after a
source law maps it into a physical generator and a clock. The conditional
transmutation formula below does not determine its running-coupling function
\(\alpha_U(P)\), cutoff or matching prescription. Their field content,
normalization and empirical inputs must be supplied separately; they cannot
be inferred from the finite gauge algebra. The comparison branches and their
limitations are stated in Section~\ref{quark-source-ew-matching}.

The universe-level capacity relation is
\[
 N=\log M_0(\mathfrak U_N),\qquad
 \mathfrak F_{r,0}(D_\star)=\{D_\star\},\qquad N=\log D_\star.
\]
Here \(M_0\) is the multiplicative size of a correctable public record code.
The finite readback is a compound-confusability independence number under a
declared coupled checkpoint family \cite{ophobservers}. Deflation,
monotonicity and greatest-fixed-point iteration require their respective
scalarization, faithful-carrier and confusability-reflection hypotheses.
A fixed-cutoff capacity calculation does not supply the universe-level
carrier or the common screen/electroweak load identification.

\section{The current response and its matter representation}

The native protected geometric object has type \(3\oplus3'\). Twelve ports
count interfaces; they do not count independent Lie generators. The tested
native exchange response closes at seventeen currents. A specified law
instead closes at
\[
 \mathfrak u(3)\oplus\mathfrak{so}(3),\qquad\dim=12.
\]
An additional coherent Hamiltonian filter has an exact invertible
control-coordinate map into this same algebra. This is a conditional control
construction, not a derivation that ordinary C/E state repair selects the
physical gauge response. The full exchange-plus-class-sum controller has
eighteen directions. Its finite filters are invertible; a fixed-bias
long-time limit has six observable directions, whereas a joint bias/update
limit with bounded drive approaches twelve. The limit and supplied baths
must be specified before making a physical current claim.

The finite formal results are implications of their typed premises. An
injective twelve-port response map has dimension twelve; it does not show
that the complete native response admits that map. The Lie-type selection
also consumes compact-simple classification data, a bound on the centre,
and an exclusion of the zero-centre alternative. The committed matter
instance supplies its spin lift, central weights and exterior mask. The
family theorem compares a finite band table within its declared window.
Checking these implications does not independently select their premises.
In particular, observability of a relative central phase is not by itself a
proof that the carrier action fixes the whole centre of the same complete
response algebra.

\subsection{Carrier covariance and physical rotations}

The specified current construction uses the port vectors in its colour
\(3\), while the family construction uses an equivalent carrier triplet
on another factor. The real part of the colour generators implements
rotations about those port axes. The carrier recharting premise acts by
inner automorphisms on the current algebra. These facts establish carrier
covariance; they do not establish which transformations observers recover
as physical spatial rotations. Equivalent finite-group representations on
different factors need not be one physical action. Conversely, declaring
the same action to be both physical frame rotation and colour transport
does not establish their required independence.

The existing spin--matter join tests a locked transport and then separates
the factors with an explicitly chosen inverse internal transport. This is
an additional construction choice. A physical identification must exhibit
the spatial action, the internal action and their commutation on the same
record/field domain, or calculate the observable consequences of a residual
lock. The distinct \(3'\) character alone does not settle the role of the
colour \(3\). Likewise, a rank-three local carrier frame is not a proof of
three-dimensional physical propagation on a federation. Neither a physical
symmetry violation nor its resolution follows solely from the finite
representation match.

\subsection{Matter lift and interaction channels}

The global finite gauge identification is
\((\SU(3)\times\SU(2)\times\U(1))/\mathbb Z_6\) under its stated faithful
representation and central-kernel hypotheses. It fixes allowed
representations, not their excitation energies. The following exact lift
and interaction calculation connects the chosen current operators to the
matter channels rather than relying on an abstract dimension match.

\section{Matter and interaction readout from the finite response}
\label{sec:repair-matter-readout}

The finite response determines more than the abstract gauge Lie type. Its
matrix generators support an exact, conditional matter decoder and local
interaction tensors. A representation is distinct from a physical particle:
a particle identification also consumes its spin, statistics, propagator,
interaction and observer-clock interpretation.

\subsection{From response operators to a fifteen-state matter module}

Let a response generator on the six-dimensional carrier have blocks
\(D=\operatorname{diag}(A,B)\), with \(A\in\mathfrak u(3)\) and
\(B\in\mathfrak{so}(3)\). Write \(B=\widehat w\). The declared spin lift and
central trace redistribution give
\[
 L(D)=\operatorname{diag}\left(
 A,-\frac{i}{2}w\cdot\sigma-\frac{\operatorname{tr}A}{2}I_2\right).
\]
This is a Lie-algebra homomorphism. The complete twelve-direction response
and all sixty-six bracket pairs are checked exactly. The central direction,
after clearing denominators and choosing the weak-positive orientation, has
primitive integer character
\[
 q=\operatorname{diag}(-2,-2,-2,3,3).
\]
The lift and trace redistribution are hypotheses. A native weak rotation
through \(2\pi\) is the identity on its vector carrier and acts as minus the
identity on the lifted doublet. Thus the lifted module is not a
subrepresentation of the original six-dimensional response. Global descent
and a physical current interpretation require the corresponding additional
identifications.

\begin{proposition}[Conditional matter decoder]
For the lifted five-dimensional one-body action, the representation
\[
 M=\Lambda^2\mathbb C^5\oplus\Lambda^4\mathbb C^5
\]
has dimension fifteen. Its charge and quadratic Casimir operators identify
the following irreducible components without a field-name input table.
\end{proposition}

\begin{center}
\begin{tabular}{ccccc}
\toprule
Component & Color & Weak dimension & \(Y=q/6\) & Dimension\\
\midrule
\(Q\) & \(3\) & \(2\) & \(1/6\) & \(6\)\\
\(u^c\) & \(\overline3\) & \(1\) & \(-2/3\) & \(3\)\\
\(d^c\) & \(\overline3\) & \(1\) & \(1/3\) & \(3\)\\
\(L\) & \(1\) & \(2\) & \(-1/2\) & \(2\)\\
\(e^c\) & \(1\) & \(1\) & \(1\) & \(1\)\\
\bottomrule
\end{tabular}
\end{center}

\begin{proof}
Split \(\mathbb C^5=V_c\oplus V_w\), of dimensions three and two and
charges \(-2\) and \(3\). Then
\[
\Lambda^2(V_c\oplus V_w)=\Lambda^2V_c\oplus(V_c\otimes V_w)
 \oplus\Lambda^2V_w,
\]
with integer charges \(-4,1,6\). The degree-four summands are
\(\Lambda^3V_c\otimes V_w\) and
\(\Lambda^2V_c\otimes\Lambda^2V_w\), of charges \(-3,2\).
Color and weak Casimirs separate the representations, and Cartan weights
distinguish the conjugate color character. Exact projectors give the same
answer after a rational change of basis mixing unlike charge sectors.
\end{proof}

The auxiliary exterior construction has two additional invariant lines, the
vacuum and top exterior degree. Removing them is part of the declared
fifteen-state module. Charge traces alone do not remove neutral invariant
states. Its even auxiliary parity also does not establish physical fermion
statistics. Physical two-component Weyl fields require an independent CAR
algebra and spin interpretation. The anomaly-free charge table and common
central kernel are conditional representation results; they do not identify
laboratory species or select a generation count.

\subsection{One Higgs doublet and the exact interaction tensors}

For two same-chirality physical Weyl fields, the antisymmetric spinor
contraction makes the internal coefficient symmetric. The Lorentz-scalar
internal pair space is consequently \(\operatorname{Sym}^2M\), of dimension
120. A color-Casimir kernel and weak-Casimir projector identify one doublet
of hypercharge \(-1/2\) and two of hypercharge \(+1/2\). Their constituent
integer-charge pairs are \((-4,1)\), \((-3,6)\), and \((1,2)\).
These are interpolating channels; their existence supplies no bound-state
pole, condensate or elementary-boson commutation relation.

\begin{proposition}[Conditional one-doublet interaction]
Declare one scalar doublet \(H\) of hypercharge \(1/2\). Each of the three
channels above has a one-dimensional invariant contraction space. They give
\[
 Q u^c H,\qquad Q d^c H^*,\qquad L e^c H^*.
\]
The same scalar doublet and its conjugate occur in all three expressions.
\end{proposition}

\begin{proof}
Restrict the exact pair generators to each color-singlet doublet. The tensor
sum of that doublet generator and the generator of \(H\), or its conjugate,
has a one-dimensional common kernel. Equivalently, its coefficient matrices
\(C_\alpha\) satisfy, for each of the twelve source currents,
\[
 T^t C_\alpha+C_\alpha T+
 \sum_\beta(T_H)_{\beta\alpha}C_\beta=0.
\]
Here the conjugate action is \(-T_H^t\) in the skew-Hermitian convention.
The physical pairing kernel \(\epsilon\otimes C_\alpha\) is antisymmetric.
The identity \(S^t\epsilon S=\epsilon\) for \(S\in\mathrm{SL}(2,\mathbb C)\)
establishes its Lorentz scalar character. The executable checks the
rotation and boost generators as well as all current identities.
\end{proof}

For three declared families the coefficients are three complex
\(3\times3\) matrices, with 27 complex entries before flavor-basis
redefinitions. The pairing convention is
\[
 \mathcal L_{\rm pair}=\frac12\epsilon_{AB}
       \psi_A^tK(H)\psi_B+\mathrm{h.c.}
\]
up to an overall sign. With canonical kinetic normalization and
\(H=(0,v)\), its singular values are tree-level mass parameters;
\(v=v_{\rm SM}/\sqrt2\) in the conventional electroweak normalization.
Exact matrix-unit tests cover arbitrary complex family coefficients.
Their values, the Higgs potential and its vacuum are not selected by these
identities. The optional invariant \((LH)(LH)/\Lambda\) has six complex
symmetric family coefficients and gives \(M_\nu=\kappa v^2/\Lambda\).
It is an additional dimension-five effective interaction, with no selected
coefficient, scale or ultraviolet completion.

If both family factors are identified with the same unbroken irreducible
rank-three source representation and the Higgs is family-neutral, the exact
commutant has dimension one. In that branch each Yukawa intertwiner is
\(yI_3\), giving equal singular values. A nonscalar background can split the
band; choosing that background is an additional physical hypothesis. The
source family-band theorem therefore does not imply the observed hierarchy.

\subsection{Scope of the readout}

A declared coherent filter on the native six-dimensional response admits an
exact invertible twelve-control map into the same compact algebra; the
spin/exterior and Higgs tensor maps above apply after this change of controls.
It must not be identified with state cooling. For a rank-three projector
\(P\), the cooling channel
\(\Phi(\rho)=P\rho P+\operatorname{tr}((I-P)\rho)P/3\)
has dual \(\Phi^*(B)=PBP+\operatorname{tr}(PB)(I-P)/3\), which is not the
Hamiltonian map \(h\mapsto PhP\). Its classical recorded branches lose the
relative phase of the two triplets.

The executable calculations reconstruct the response, representation
projectors and interaction tensors from repository-local source operators,
with exact rational or algebraic identities separated from numerical
residuals. They use no measured particle masses. They supply neither a
physical chiral field net nor a source-selected Higgs vacuum, Yukawa matrix
or mass pole. A relaxation gap, finite cavity frequency or commuting
internal modular offset requires its own mass-shell identification before
it can be called a physical rest mass.


\subsection{Kinetic normalization and physical conventions}

For a four-component Dirac convention with positive kinetic metrics,
\[
 Y_{s,\mathrm{can}}=
 Z_L^{-1/2}Y_s Z_R^{-1/2}/\sqrt{Z_H}.
\]
A literal all-left Weyl bilinear instead has a transpose on the left
normalization factor. Conversion between these conventions must precede
comparison of coefficients. The up and down sectors share \(Z_Q\) because
their left fields belong to the same weak doublet.

A record matrix can change a coordinate without changing the mass. For
example, \(Z_L=Z_R=R>0\), \(Y_s=g_sR\) gives
\[
 Y_{s,\mathrm{can}}=g_sI/\sqrt{Z_H},\qquad
 \det(pR-hg_sR)=\det(R)(p-hg_s)^3.
\]
Thus a hierarchical Gram matrix by itself need not generate a hierarchical
physical vertex. Eliminating an invertible heavy block in a source kernel
also changes its Higgs derivative:
\[
 \partial_H(A-BD^{-1}C)=A_H-B_HD^{-1}C-BD^{-1}C_H
                       +BD^{-1}D_HD^{-1}C.
\]
Both the kinetic and Higgs-dependent parts must come from the same source.

Species are identified by the represented charges, Casimirs and interaction
channels, with conjugate-charge and chirality conventions stated. Within a
species, sorting a nondegenerate singular spectrum is a permissible naming
convention for mass eigenstates; it is not a derivation of three families
or their weak-current couplings. For canonical Dirac matrices, write
\(U_{sL}^{\dagger}Y_sU_{sR}=\operatorname{diag}(y_s)\). The charged-current
matrix is \(V=U_{uL}^{\dagger}U_{dL}\), up to the corresponding ordering and
rephasing conventions. Generation need not be a conserved quantum number.
A physical CP test uses a rephasing-invariant quantity such as the
Jarlskog invariant, not the presence of a complex matrix entry
\cite{pwjarlskog1985commutator,pdg2024ckm}.

In the executable all-left convention, a directed complex family Yukawa
appears as \(\overline Y\), with the reverse block its transpose. The
creation coefficient is \(\epsilon\otimes\overline K\). Real symmetric
tests cannot detect a mistaken conjugation, so the finite checks include
complex matrix units and all branches of the 180-dimensional Nambu
representation. Its supplied spatial Weyl/Higgs background is distinct from
a source-derived interacting field.

\subsection{Reusable finite kinematic and interaction identities}

The strengthened finite family interface selects rank three under its
single-object, faithful-action and operational-cost hypotheses. A separate
conditional argument from intrinsic CP capability and weak-sector UV
completeness gives \(3\le N_g\le5\). Neither statement attaches a physical
family field or excludes extra light sectors. The ten-row exterior census
follows from the 32 subsets of five modes: bidegrees \((c,w)\) have
multiplicity \(\binom3c\binom2w\) and charge \(-2c+3w\). Removing vacuum
and top degree gives ten bidegrees. A supplied anomaly-free chiral mask and
central weights recover the same six-element kernel. Partition readout can
transport this dictionary under declared maps; it does not select the maps,
physical parity or laboratory charges.

With a declared scalar quartic, breaking scale and gauge couplings, the
finite electroweak-breaking theorem identifies the vacuum shell, its
one-dimensional stabilizer, three gauge-orbit Hessian zero directions and
the tree gauge-mass ratio. The slot dictionaries then give photon charge
\(eQ\), neutral shape \(g_Z(T_3-Q\sin^2\theta_W)\), and charged shifts
\(g_W/\sqrt2\) on doublets, zero on singlets. These are identities of the
declared infinitesimal action and pairing, with supplied constants. They
neither construct physical current operators nor supply mixing, decay
amplitudes or a Fermi limit.

On a declared Lorentz module and mass shell, free plane waves persist under
the declared translation flow, and their phase along their own frame has
proper-time derivative \(m\). For \(m>0\) the internal-clock periods are
\((2\pi/m)\mathbb Z\); the zero-mass clock is constant and has no
fundamental period. This is conditional kinematics, without a
source-selected mass or clock. Likewise, for a declared complex algebra
with \([X_i,X_j]=iq\epsilon_{ijk}B_k\), the Pauli identity is
\[
 (\sigma\cdot X)^2=(X\cdot X)I-q\sigma\cdot B.
\]
In the stated orbital/spin normalization it reads as \(g=2\). The identity
does not construct a gauge potential, electron, anomalous moment or
physical coupling.

\section{Uncalibrated response laws and their exclusions}

\subsection{Complete memory records and normalized feedback}

The complete twelve-state memory uses the native events with probabilities
\(p=1/12\), \(q=1/6\), a specified reset state and declared rotations. All
128 native raw branches are combined with all twelve old memory states;
conditioning on successful events is not substituted for this 1,536-branch
law. The stationary projector average
\[
 Y=\sum_k\pi_kP_k\succeq0,\qquad\operatorname{tr}Y=1
\]
has coefficient ratios \(1:195.11703757:3579.82624571\). A two-angle
record inverse cancels the finite-pulse bias exactly. The resulting numerical
allowance is conditional on its floating-point and mixing estimates; it is
not an outward-rounded, finite-shot or physical uncertainty certificate.

Gating this memory with the complete thirteen-outcome matter instrument and
both components of each Higgs tensor gives different sector coefficients.
The entire declared geometric catalog \(t\in\{\pi/5,\pi/3,\pi/2\}\) fails
the charged hierarchy comparison. Its zero coefficients are retained.
A separate finite-angle inverse reconstructs a declared static mean-feedback
Hamiltonian from fixed-coordinate preparation and detector probabilities.
Averaging the memory and executing a correlated quantum/classical history
are different operations; neither calculation proves their dynamical
equivalence.

\begin{proposition}[Normalized-feedback scale bound]
For a probability average of rank-one projectors on a three-dimensional family space,
\(Y\succeq0\), \(\operatorname{tr}Y=1\), and
\(1/3\le\lambda_{\max}(Y)\le1\). With common canonical kinetics,
unit-modulus neutral Higgs tensors and a common Higgs norm \(v\), the three
family action masses in each sector sum to \(v\). Any two sector maxima
have ratio at most three.
\end{proposition}
\begin{proof}
The nonnegative eigenvalues sum to one, so their largest lies between their
mean and their sum. The specified Higgs tensors multiply each family
singular value by \(v\). Apply the bounds to both sectors.
\end{proof}

This is a sector-normalization obstruction under the common-coupling and
canonical-kinetic hypotheses. It does not bound a sector's smallest
eigenvalue away from zero and therefore does not exclude arbitrary
within-sector hierarchies. A sector-dependent overall strength or a
different physical kinetic response changes those hypotheses and needs its
own derivation. CP and mixing are separate questions: Yukawa matrices
simultaneously real in one weak basis have zero CP-odd invariant, while a
common matrix up to sector scalars has aligned left singular frames and
trivial charged-current mixing for a simple spectrum. Different real
matrices can mix, and general complex positive trace-one matrices need
not have vanishing CP violation.

Unconditional event-duty pulses change this normalization. They must use
the complete correlated record law. In particular, for reset projector
\(P_2\), the C-duty response is
\[
 Z_C=Y-(1-p)P_2,
\]
rather than \(pY\). The always/C/E/CE policies have traces
\(1,p,q,pq\). Their sector-maximum spans in the declared implementation are
approximately \(1.0233,2.0572,1.1528,2.3695\); the E and CE responses have
rank two. These calculations exclude those policies as the charged mass
prescription. They do not exclude every repair law on the carrier.

\subsection{Quadratic response and the affine-family obstruction}

For the native quadratic comparison, put
\[
 a=\log(25/8),\quad b=\log(8/5),\quad J=a-b,\quad\Omega=a+b.
\]
The three eigenvalues of its family matrix \(F\) are
\[
 \frac{61}{450},\qquad\frac{41}{225}\pm\frac{7\sqrt5}{750}.
\]
The direct positive prescription has ratios approximately
\(1:1.1903:1.4982\), far from a charged hierarchy. Allowing arbitrary complex
affine coefficients does not repair this class.

\begin{proposition}[Affine quadratic-family exclusion]
Let \(Y=uI+wF\), \(u,w\in\mathbb C\), and let
\(m_{\min}\le m_{\mathrm{mid}}\le m_{\max}\) be its singular values.
With \(\varphi=(1+\sqrt5)/2\),
\[
 m_{\max}\le\varphi^2m_{\mathrm{mid}}+\varphi m_{\min}.
\]
Consequently this class cannot reproduce the charged-lepton comparison
windows used here.
\end{proposition}
\begin{proof}
For ordered eigenvalues \(f_1<f_2<f_3\) of \(F\), the complex points
\(z_i=u+wf_i\) are collinear. Their adjacent spacings have ratio
\(\varphi^{-1}\). Solving the affine relation for each endpoint and applying
the triangle inequality bounds either endpoint modulus by
\(\varphi^2\) times either of the other moduli plus \(\varphi\) times the
remaining modulus. The middle point is a convex combination of the
endpoints. Applying these bounds to the largest modulus proves the stated
sorted bound. Conditional on \(m_{\mathrm{mid}}/m_{\min}\le207\), it gives
\(m_{\max}/m_{\min}\le543.551\), whereas the broad pole-ratio comparison
requires at least \(3470\). This numerical bound is not uniform over all
affine coefficients without the middle-to-light ratio constraint.
\end{proof}

This result concerns a fixed quadratic action and its affine family. Its
anticommutator pole locations are state-independent; an occupation-energy
versus fermionic-gap factor of two, a common clock rescaling or a basis
change cannot change their ratios. Interactions that genuinely change the
kernel require a separate calculation.

The pole-ratio windows \([206,207]\) and \([3470,3485]\) are only a
diagnostic target for these dimensionless tree coefficients until a
physical matching map is supplied. For example, the adopted
\(\overline{\mathrm{MS}}\) comparison at \(M_Z\) has lepton ratios about
\(210.7\) and \(3578\), rather than the pole values about \(206.77\)
and \(3477.4\) \cite{antusch2025running}. The affine inequality still
excludes that comparison, but its numerical bound must use that scheme's
middle-to-light ratio. Ratios do not remove all scale and scheme dependence.

Other explicit laws also fail: a five-memory law gives ratios near
\(1:30.84:452.62\); a two-rotation law gives
\(1:382.23:5909.43\); the admitted auxiliary-dressing class bounds its two
smallest subthreshold lepton poles by a ratio below 32. These are specified
model exclusions. A near numerical ratio obtained by searching many
prescriptions is not a blind prediction; the selection ancestry and all
admitted alternatives must be retained.

\subsection{Full repair and physical propagation}

An exact decision experiment first tests a local defect candidate in the
unchanged total-load-two native C/E domain. All 720 raw outcome maps are
retained on its 78 configurations. Twelve transient double-occupation
classes feed one closed irreducible binary class of dimension 66. A scalar
configuration invariant preserved by every admitted native branch is therefore
constant. Under the separately specified independent uniform-seam/fair-coin
schedule, a double load survives \(n\) attempts with probability
\((5/6)^n\), has mean lifetime six attempt counts, and splits rather than
moving as a protected defect. In the uniform binary stationary law the
nonzero-current probability is \(5/33\), while every event has zero drop of
\(\Phi=\tfrac12\sum_p n_p^2\). Edge-disagreement costs can both rise and fall.
Thus motion and repair activity are not that potential decrease. Removing
the five seams at one port creates closed classes of sizes 11 and 55 only by
changing the allowed schedule. These are configuration-sector results, not
an exclusion of quantum, history or enlarged-carrier defects.

For this same collision candidate the full quantum instrument gives the
stronger exact identity
\[
 \sum_h K_h^\dagger P_{\rm coll}K_h=\frac56P_{\rm coll}
\]
for either C or E, including zero-probability outcomes in the retained menu.
There is no binary-to-collision feed. With the stated independent uniform
attempt schedule, arbitrary inputs entangled with an untouched reference
therefore have collision weight \((5/6)^n\) times their initial weight.
No persistent quantum subspace is supported wholly on these collisions under
that unread evolution. This does not classify other quantum or history sectors.

The actively recovered matching code has a different contract. Without
recovery, the uniform-seam C/E maps take its normalized code projector outside
the code with probabilities \(31/108\) and \(55/108\), respectively. With
the declared recorded recovery, its logical state and references are preserved
while the complete record law is independent of the logical input. A fixed
cost of those records alone cannot split logical-state masses. This does not
exclude a common energy offset. The fixed matching table, source-derived
equivariant code and a new nonscalar correlation observable cannot be silently
identified. The native invariant, localized excitation and common energy/clock
attachment remain explicit obligations of any defect route.

Passive stationary fixed-reference schedules have an exact rank-two Schur
response obstruction, including eight zero crossing weights among the
22 seams. Positive whitening preserves the null direction. The uncorrected
recorded 66-state channel has no invariant fifteen-state work code in the
tested class. Its ten-mode response conditioned on an incident event has
operator norm \(\lambda_C=7/10\) or \(\lambda_E=2/5\). For a unitary
\(U\) and any linear map \(J\),
\[
 \|UJ-JR\|_{\rm HS}\ge(1-\lambda)\|J\|_{\rm HS}.
\]
Thus it has no nonzero exact direct intertwiner to the unmeasured closed
quadratic Weyl evolution. The theorem does not exclude a field coupled to
detectors, an environment or interactions; a complete dephasing-channel
control realizes the same uniform contraction. Such a bridge must match the
complete record and continuation law. Actively corrected, copied,
controlled and refinement-dependent constructions have separate domains.

Collective charge operators can be reconstructed from shared records without
being a physical fermion field. The free spatial implementation has
conditional Weyl--Higgs dispersion on a supplied background; its cubic
staggered continuum contains two Dirac tastes and both chiralities. A
restricted momentum comparison does not prove a unique physical continuum.
Likewise, a free wedge modular/stress calculation does not prove physical
entropy stationarity or the interacting Einstein construction
\cite{ophEinstein}. The source program must connect these separate
calculations on one admissible field and observer history.

Response diagnostics now separate supplied-action spectroscopy from actual
record reconstruction. The finite interacting CAR response determines its
inverse and frequency derivative from all Lehmann residues, with canonical
sum-rule and positivity admission. Kinetic and mass forms transform together
under a change of interpolators; raw matrix eigenvalues are not invariant
mass data. A separate charge-resolved spectral census of the supplied
180-mode Weyl/Nambu symbol retains signed and dark modes, unresolved frequency
lists and numerical rather than certified exact residue ranks.

For one existing charged \(L/e^c\) pair, a finite neutral-reference experiment
reconstructs the complex two-mode transfer from complete detector/clock
records. Charge-compatible source pulses prepare a fixed-total-charge-zero
reference, and all measured continuation branches remain. On a supplied
principal logarithm window and canonical quadratic response class, the inverse
recovers the generator, residues and joint z-axis momentum and neutral-Higgs
multiplier derivatives. The retained 304 contexts have 4,864 terminal outcomes.
Withheld settings and a second duration give generator residuals at most about
\(6.40\times10^{-8}\) per unit of reconstructed clock coordinate and
complete-output distances about
\(4.36\times10^{-8}\). Both durations repeat the same elementary source word;
their composition defect is about \(2.23\times10^{-16}\). These are numerical
diagnostics, not physical error certificates. Phase-sensitive controls detect
opposite coupling signs hidden by Rabi marginals, and withheld momenta detect
an otherwise indistinguishable two-hop kinetic control. The alias-window prior
is not removed by the two durations.

This reconstructs a component of the existing supplied action. It selects no
mass coefficient, introduces no new particle candidate, supplies no dynamical
Higgs/current vertex and does not establish an interacting continuum. Its
completed coherent-control services and clock remain premises; native attempt,
abort and export histories are not replaced by unit success probability.
The executable entry points and complete readout contracts are recorded in
\path{docs/research/OPH_FIELD_RECONSTRUCTION.md}.

There is already a finite abelian quantum link with an electric--matter
interaction and an exact operator-Gauss sector; it is a restricted
four-state code, not the full nonabelian particle theory. The Euclidean
overlap calculation on the retained fifteen-state representation goes
beyond a free symbol: it checks the modified chiral projector and gauge
covariance on a nonconstant nonabelian background, with a stated
small-plaquette locality domain. Its native kernel queries retain success
and failure branches. These results leave the global chiral measure,
reflection positivity or another justified Lorentzian construction,
dynamical gauge/Higgs fields and common physical clock open. The naive
lattice's eight balanced Weyl nodes are a negative control; deleting them
does not prove decoupling. The precise retained domains are documented in
\path{model/common_source/DERIVATION.md} and
\path{model/common_source/chiral_candidate/REGULATOR_ROUTES.md}.

\section{Flavor constraints and comparison constructions}

\subsection{What a family band determines}

A source-executed protected rank-three band is a candidate family carrier.
Its identification with three physical generations is additional. If both
family factors have the same unbroken irreducible action and the Higgs is
family-neutral, Schur's lemma forces \(Y_s=y_sI_3\), hence degeneracy.
Splitting requires a physical symmetry-breaking response. A declared
transport kernel, its centered generator, overlap lines, gaps and cycle
phases are useful mathematical data only after their common source and
physical channel are specified.

For a positive Hermitian circulant square-root response
\(C=aI+bR+\overline bR^2\), \(R^3=I\), let \(b=ar e^{i\delta}\).
If its eigenvalues are the nonnegative square roots of the masses, then
\[
 Q:=\frac{\sum_i m_i}{(\sum_i\sqrt{m_i})^2}
   =\frac13+\frac23r^2.
\]
Indeed, \(\operatorname{tr}C=3a\) and
\(\operatorname{tr}C^2=3a^2+6|b|^2\). Thus \(Q=2/3\) exactly when
\(r=1/\sqrt2\). This is the Koide/circulant identity
\cite{koide1983,brannen2006}; it does not select a positive branch, the
phase, an absolute scale or a physical response. A supplied GNS balance
condition is a premise rather than a consequence of this trace identity.
Its retained equal-weight block realization uses an asymmetric choice of
blocks abstracted from the measured charged triple. Treating those blocks
alike gives a different amplitude and \(Q=1\). Thus the construction
realizes the empirical balance condition; it does not independently select
Koide's value or the fitted phase \(2/9\).

Orientation reversal can send a candidate cycle phase
\(\delta\mapsto-\delta\). Identifying that operation with physical
antiunitary time reversal requires its own intertwiner. A classical cycle
affinity is a log probability ratio, not a coherent \(\U(1)\) phase. Detailed
balance of one declared kernel does not rule out phases in another coherent
sector or in a controlled extension. A phase inference must use the actual
source process and an observable phase readout.

\subsection{Quark moduli and common-scale comparison}
\label{quark-family-strongest-overlap-branch-primitives}
\label{quark-family-boundary}

Conditional on the registered source signature containing no typed datum
that separates the stated centered-spread orbit, a generic quark package
does not define a unique quark mass spectrum. Granting an ordered profile
shape, write its centered logarithms as \(\sigma_u\widehat\ell_u\) and
\(\sigma_d\widehat\ell_d\). Independent positive changes of
\((\sigma_u,\sigma_d)\) preserve the certified gauge representations,
anomaly traces and invariant Higgs lines. Those certificates therefore
leave at least an \((\mathbb R_{>0})^2\) fiber in this granted subfamily.
Without the shape hypothesis, sector means and shapes are also undetermined.

Representative independence does not select the datum: any constant map
on a selected bridge fiber descends to its quotient, whatever its constant
value. A target-derived spread retains that ancestry after descent.
For the minimal light branch
\[
 y_u=c_u\varepsilon^6,\quad y_d=c_d\varepsilon^6,\quad
 \varepsilon=1/6,\quad
 \Delta_{ud}=\tfrac16\log(c_d/c_u)=t_1/5,
\]
the conditional bridge fixes a ratio, not either independent spread.
The candidate reciprocal-ray law fails on common-scale Standard Model
\(\overline{\mathrm{MS}}\) Yukawas at \(M_Z\)
\cite{antusch2025running}: its required product of adjacent log-gap ratios
is one, whereas the adopted comparison gives approximately \(0.83532\).
That is a rejection on the stated comparison surface, not a global
classification of quark sources.

The one-doublet tensors can type a relation between down-quark and
charged-lepton channels; they do not require equal coefficients. Under the
additional alphabet, measure-balance and register-faithfulness hypotheses,
finite enumeration selects \(\{1/3,1,3\}\). Ranking its six generation
assignments against measurements is target-informed. The selected
assignment fails the adopted joint FLAG comparison \cite{flag2024}.
The assignment \(y_b=y_\tau\), \(y_s=y_\mu/3\), \(y_d=3y_e\)
is the Georgi--Jarlskog texture \cite{pwgeorgi1979new}; its historical
ancestry and target-informed ranking are distinct from the conditional
alphabet enumeration. The retained six-quark mixed-convention table also
comes from a target-conditioned template search and is not an independent
common-scale prediction.
Similarly, directly identifying two real \(A_5\) residual axes with the
Cabibbo angle is excluded in that class: the least nonzero angle is about
\(20.905^\circ\), with sine \(0.356822\), rather than the measured
approximately \(0.225\). This does not exclude general complex flavor
transport or additional dynamical misalignment.

\subsection{Conditional fixed-pixel and electroweak matching charts}
\label{quark-conditional-fixed-p}\label{quark-source-ew-matching}

Explicit fixed-\(P\) numerical constructions evaluate running-chart
coordinates after their flavor law, absolute clock and matching prescriptions
are frozen. They are conditional constructions: their discrete choices and
some attachments have target ancestry. Freezing them does not erase that
ancestry. GeV-valued textures assembled from
mixed-scale quark rows are not dimensionless Yukawa matrices at a common
scale. A physical comparison must transport all six coordinates, including
the top extraction convention, to one scheme and scale, emit the running
Higgs expectation value in that scheme, and then use
\(y_q(\mu)=\sqrt2m_q(\mu)/v_{\rm SM}(\mu)\).

The same boundary applies to an electroweak matching stage that keeps only
the top Yukawa in a Standard Model truncation. A root of that declared map
is not a unique source-selected vacuum. Conditional solutions of
\(\lambda=\beta_\lambda=0\) select a critical boundary after its
couplings, truncation, clock and matching are supplied. Positive-capacity
top bands and fixed-pixel near-agreements do not remove those premises.

\paragraph{Running content and branch provenance.}
The historical transmutation recipe used the MSSM one-loop coefficients
\((33/5,1,-3)\), with \(g_1=\sqrt{5/3}\,g_Y\). The supplied Standard Model
content of three fifteen-state generations and one scalar Higgs doublet
instead gives \((41/10,-19/6,-7)\) in that convention, or
\((41/6,-19/6,-7)\) for \((g_Y,g_2,g_3)\), with
\(\mu\,dg_i/d\mu=b_i g_i^3/(16\pi^2)\). No superpartner threshold content
is derived here. The MSSM packet is therefore an additional benchmark
recipe, not the running of the stated matter construction. Correcting it
requires a common normalization, threshold/matching specification and
recalculation of the closure and transmutation maps; it need not preserve
a unification root or the previous numerical agreement.

The near-measurement archived W/Z tuple used \(P=1.63094\). It must not
inherit the label of the distinct strict closure branches near
\(1.6309682\) and \(1.6309721\). Values from those branches, and fitted
W/Z reparametrizations, have separate provenance. The legacy absolute
normalization \(E_\star\simeq1.22089\times10^{19}\,\mathrm{GeV}\) is the
Planck energy \(\sqrt{\hbar c^5/G}\) using measured \(G\); rewriting it
as a clock coordinate does not derive that input. The retained source
records and their current classifications are indexed in
\path{code/particles/MASS_CANDIDATE_STATUS.md}.

\paragraph{Criticality is a conditional boundary.}
The Higgs/top calculation imports Standard Model running and empirical
gauge inputs \cite{buttazzo2013}. The scale
\(E_\star e^{-\pi}P^{-1/6}\) was proposed after exposure to the implied
boundary scale; it is a post-exposure candidate, not an independently
selected source boundary. The retained comparison with its supplied
\(172.1\,\mathrm{GeV}\) top reference gives a two-loop curve coordinate
\(124.7256\,\mathrm{GeV}\) for the Higgs and an implied boundary scale
\(2.9545\times10^{17}\,\mathrm{GeV}\). The older \(125.72\,\mathrm{GeV}\)
curve display does not describe that comparison. The one-loop top range
is \([163.0899,170.7539]\,\mathrm{GeV}\), so it does not bracket the
adopted top reference. These are arithmetic comparisons of declared
coordinates, not physical confidence intervals: the Higgs tree coordinate,
QCD-converted top and experimental mass conventions still require common
matching and theory uncertainties. See
\path{code/particles/runs/calibration/d11_criticality_comparison.json} and
\path{code/particles/runs/calibration/d11_boundary_scale_selection_audit.json}.

\paragraph{Strong-CP boundary.}\label{quark-family-strong-cp-boundary}
Strong CP is independent of the spread and normalization obstructions.
The source does not emit the anomaly-invariant \(\bar\theta\) or prove
that it vanishes. That requires a common-scale physical quark determinant
phase and a source law for the QCD topological angle; GeV-valued mass
textures supply neither.

\subsection{Charged shape, determinant and finite-register candidates}

On the declared ordered carrier \((-1,x_2,1)\), with
\(x_2=-0.5175863354681689\), two supplied source-support scalars
\((\eta,\sigma)\) determine the centered charged log triple
\[
 \ell_e=-\frac{(3+x_2)\sigma-\eta}{6},\quad
 \ell_\mu=\frac{x_2\sigma-\eta}{3},\quad
 \ell_\tau=\frac{(3-x_2)\sigma+\eta}{6}.
\]
Its sum is zero. The downstream algebraic readout is
\(m_i=g_e e^{\ell_i}\); neither the displayed carrier coordinate nor the
physical source/charged attachment follows from this formula.

\begin{theorem}[Charged absolute-scale underdetermination]
Any readout factoring only through this centered log triple and its
difference invariants leaves \(g_e>0\) free.
\end{theorem}
\begin{proof}
For every \(c\in\mathbb R\),
\(Y_e(c)=e^c\operatorname{diag}(e^{\ell_i})\) has the same centered logs
and differences. Its determinant is \(e^{3c}\), and its masses scale by
\(e^c\). Thus the stated data determine a common-shift quotient, without
a section selecting the determinant scale.
\end{proof}

An uncentered physical response could supply the affine-covariant anchor
\[
 A_{\rm ch}(\ell+c\mathbf1)=A_{\rm ch}(\ell)+c,\qquad
 A_{\rm ch}=\tfrac13\log|\det M_e|.
\]
The common-shift theorem does not exclude such an additional datum. A formal
source determinant character must be attached to the physical charged
determinant line before this identity becomes a mass law.

\begin{proposition}[Twenty-four-register rate no-go]
An oriented register with twenty-four slots, a twenty-four-stage automaton, or
a decomposition into twenty-four formal update steps determines no physical
frequency, Hamiltonian gap, or mass scale. Indeed, rescaling every physical
generator by \(H\mapsto\lambda H\), \(\lambda>0\), preserves the register,
transition graph, grading, and settled quotient while multiplying all
generator gaps and rates by \(\lambda\). A repair iteration counter is likewise
not an operational clock. A physical rate requires a source-derived clock
instrument, calibrated comparison line, and residual bound.
\end{proposition}

\begin{theorem}[Conditional charged determinant-clock attachment]
Let \(L_{24}\) be a one-dimensional source-derived operational clock line
with a distinguished calibrated positive gap vector \(\ell_{24}\) of norm
\(\Delta_{24}\). Suppose there is a
refinement-natural, quotient-visible norm-preserving line isomorphism
\[
\Theta_e:L_{24}^{\otimes3}\longrightarrow\det M_e
\]
to the physical charged-lepton determinant line, together with the load-bearing
element compatibility
\[
\Theta_e(\ell_{24}^{\otimes3})=\det(M_e).
\]
Then the common charged
mass scale is fixed by
\[
g_e:=|\det M_e|^{1/3}=\Delta_{24}.
\]
More generally, if the element-compatible map rescales norms by a separately
derived factor \(\kappa>0\), then \(g_e=\kappa^{1/3}\Delta_{24}\). Thus the
implication from a calibrated clock-line
attachment to the determinant scale is exact.
\end{theorem}

\begin{proof}
The distinguished tensor-cube vector has norm
\(\|\ell_{24}^{\otimes3}\|=\Delta_{24}^3\). Element compatibility and norm
preservation give \(|\det M_e|=\Delta_{24}^3\); take the positive cube
root.
\end{proof}


These hypotheses require a source-derived clock line and an element-compatible
physical attachment. Counting slots or writing a norm-preserving line isomorphism
without the distinguished-vector equality does not construct either.

\begin{proposition}[Finite charged-register realization]
\label{prop:charged-cfq-digital-model}
The stipulated charged register-and-path class is nonempty. There is an explicit
finite algebraic model with register dimensions
\[
(50,31,10,512,77,21,27,5),
\]
connected transition graphs, normalized tracial states, rank-one event
projectors, and the eight declared path weights
\[
\left(\frac1{50},-\frac1{31},-\frac1{310},\frac1{512},
\frac1{77},\frac1{21},\frac1{27},\frac1{135}\right).
\]
Each noncentral event in \(B(H_r)\) admits an accepted/rejected central record
dilation in \(B(H_r)\otimes D_2\). The graph family is covariant under the
sixty proper icosahedral rotations, and normalized path weights are unchanged
under inert ancillary stabilization \(X\mapsto X\otimes I_k\).
\end{proposition}

\begin{proof}
For each declared connected graph, the diagonal matrix units and both directed
units on every edge generate \(E_{ij}\) along graph paths and hence generate
the full matrix algebra. The normalized trace of a rank-one event is the
reciprocal register dimension, and tensor products multiply the two depth-two
weights. For an event \(P\), the pinching map
\[
\mathcal E_P(X)=PXP+(I-P)X(I-P)
\]
is a unital trace-preserving conditional expectation. Its instrument writes
\(P\rho P\) and \((I-P)\rho(I-P)\) into orthogonal accepted/rejected pointers in
the central \(D_2\) factor. Explicit permutation intertwiners preserve graph
incidence and ranks for all sixty proper rotations. Finally,
\(\tau_{dk}(P\otimes I_k)=\tau_d(P)\), and the partial-trace coarse map retracts
the inert embedding. These facts construct one model of the declared schema.
\end{proof}

Proposition~\ref{prop:charged-cfq-digital-model} proves formal nonemptiness.
It supplies no derivation of the charged source. The eight dimensions, path
automaton, coupling degrees, grading, signs, clock, and scalar response are
model inputs. The verifier exhausts that automaton, not the physically
admissible OPH path space. The construction resolves the abstract
event-versus-central-record issue at fixed cutoff. Physical response
selection, cofinal screen refinement, and charged-sector attachment are not
supplied.

The dimensions, paths, signs, state, scalar response and clock are declared.
This proposition proves that their finite schema is realizable; it does
not select that schema from the physical source.

\begin{proposition}[Conditional charged nature and singularity transport]
\label{prop:charged-nature-pole-transport}
Let \(M_F\) be the positive face endpoint. Suppose a physical chiral
three-family carrier supplies a natural unitary \(J_L\), canonical kinetic
metrics, and
\[
 \frac{v^2}{2}\widehat Y_e\widehat Y_e^\dagger
 =J_LM_F^2J_L^\dagger.
\]
Then its positive left charged response is \(J_LM_FJ_L^\dagger\). Suppose in
addition that the exact renormalized charged kernel is supplied, its singular
lines satisfy the declared infrared and regularity conditions, and a
finite-register/Dyson map has singularity readout
\(J_LM_FJ_L^\dagger\). Regular
invertible field changes preserve the zero set, and a regular Nielsen
factorization \(\partial_\xi K=A_\xi K+KB_\xi\) makes each simple zero gauge
independent.
\end{proposition}

\begin{proof}
The first statement follows from uniqueness of the positive square root.
Multiplication of \(K\) on the left and right by analytic invertible matrices
multiplies its determinant by nonvanishing factors. Jacobi's identity gives
\[
 \partial_\xi\det K=(\operatorname{tr}A_\xi+
 \operatorname{tr}B_\xi)\det K,
\]
so a simple zero cannot move while the Nielsen factors are regular. These are
the standard mixed-fermion and gauge-independence implications used in the
pole literature~\cite{kniehl2013,gambino1999}; the stable charged-line
infrared qualification follows the infraparticle boundary
\cite{buchholz1986,duch2023}.
\end{proof}

The displayed square identity and singularity attachment are physical
premises. The free witness \(K_0(s)=sI-M_F^2\) does not supply an interacting
charged-lepton kernel. A determinant-preserving change of three
ordered masses by \((2,1/2,1)\), with coherently changed mass, Yukawa and pole
matrices, preserves the tested algebraic relations unless the endpoint
identity \(M_F=g_{\rm end}S_F\) is separately enforced. Existence and internal
consistency are weaker than source selection and parameter identification.

\paragraph{Target-informed response architecture.}
\label{charged-mcpr-frontier}
The eight-register candidate stipulates its path weights, phase, amplitude
and determinant exponent. Its declared response map is a contraction with
constant \(1.5\times10^{-3}\), giving ratios approximately
\(m_\mu/m_e=206.7683\), \(m_\tau/m_e=3477.3655\).
A segregated comparison is about 84 ppm below the adopted charged triple.
Those architecture choices were made with knowledge of the target; a
runtime-reference-free evaluation is therefore a target-informed model
coordinate. Neither this proximity nor the Koide identity derives the
physical source law. A normalized rank-one path weight is not by itself a
\(\mathbb Z_6\) character or a determinant.

The retained inverse comparison has three target-informed spectral
coordinates: phase, modulus and scale. The absolute 84 ppm display uses
the branch \(P\simeq1.6309720958588974\) and the externally normalized
\(E_\star\); it is not the relative accuracy of the two mass ratios.
The associated ratio comparison reports residuals about \(+0.0865\) ppm
for \(m_\mu/m_e\) and \(+0.0790\) ppm for \(m_\tau/m_e\). The former
lies outside the adopted componentwise one-standard-uncertainty interval
of halfwidth about \(0.0221\) ppm. This interval check supplies neither a
joint statistical test nor a physical significance without covariance,
matching and model-error information. The comparison and formula-selection
ancestry must remain attached to the displayed proximity.

\paragraph{Electromagnetic scale inversion.}\label{charged-kappa-interval}
On a declared on-shell transport branch, a common rescaling
\(m_i\mapsto e^\kappa m_i\) gives
\[
 P_{\rm lep}(\kappa)=\frac1{3\pi}\sum_i
       [2\log(M_Z/m_i)-5/3],\qquad
 \frac{dP_{\rm lep}}{d\kappa}=-\frac2\pi.
\]
An empirical photon endpoint, target-anchored ratios, hadronic payload and
scheme bridge can therefore identify \(\kappa\). They add data excluded
from the common-shift theorem's antecedent list. The retained interval
calculation gives central masses approximately
\((0.5089,105.22,1769.5)\) MeV with comparison intervals containing the
adopted measurements. Its anchor-gap uncertainty is correlated with its
hadronic payload; treating them independently enlarges the interval.
Neither the rectangle nor the payload-coherent interval is source-only.
Their interpretation requires the declared higher-order and scheme
uncertainties, rather than a physical mass claim from centered data. The
higher-order leptonic subtraction uses the measured charged triple; solving
back for its common scale is therefore a calibration consistency check,
not an independent falsification target for those masses.

The exact one-loop spacelike Ward response offers another conditional
scale attachment. For fixed centered shape \(\ell\) and nonzero spacelike
momentum, define
\[
 I(z)=\int_0^1x(1-x)\log(1+zx(1-x))\,dx,\qquad
 \mathcal W_Q(\mu;\ell)=\frac2\pi\sum_i
             I(q^2e^{-2(\mu+\ell_i)}).
\]
Since \(I'(z)>0\), this response is strictly decreasing in \(\mu\),
with limits infinity and zero. A supplied source-complete endpoint pair
and nonleptonic subtraction therefore select a unique common scale.
They are additional inputs to the centered-data reduct. Higher-order
monotonicity, source ancestry and the clock/atomic transport factors
are separate requirements; inverting measured endpoints is calibration.

\subsection{Conditional neutrino deformation}

Fix an isotropic Majorana reference \(M_0\), with real \(a,\rho>0\) and cycle
phase \(\Omega\). Supplied same-label gap and overlap data
\(g_e>0\), \(0\le\omega_e<1\), \(e\in\{12,23,31\}\), give
\[
 q_e=\sqrt{g_e(1-\omega_e)},\quad
 \eta_e=\log q_e-\tfrac13\sum_f\log q_f,\quad
 \mu_e=\frac{e^{\eta_e}}{\frac13\sum_f e^{\eta_f}}.
\]
The gaps and overlaps in this comparison inherit a template transport kernel
and candidate line lift. Centering removes a common scale; it does not
establish their source ancestry or physical family attachment.

\begin{theorem}[Fixed isotropic reference plus same-label scalars suffice for the intrinsic deformation]
Fix the isotropic reference \(M_0\), equivalently its parameters
\((a,\rho,\Omega)\). Assume the same-label gap and overlap scalars are supplied on all realized
arrows. Then the full
intrinsic neutrino mass-eigenstate bundle factors through the scalar certificate
\[
\bigl(g_e,\omega_e\bigr)_{e\in\{12,23,31\}}
\]
or equivalently through the centered eta-class \([\eta_e]\). No additional raw matrix payload is
needed to form the intrinsic selector, the depressed-cubic spectrum, the ascending gaps, or
the intrinsic spectral subspaces. When the singular spectrum is simple and strictly positive,
canonical Takagi congruence fixes an ordered column basis up to column signs. If a simple spectrum
contains one zero singular value, its zero column retains an arbitrary \(U(1)\) phase while the
nonzero columns are fixed up to signs. At a degeneracy, including the
isotropic \(s_0=s_1\) point, only the degenerate spectral subspace and its projector are fixed.
Ascending singular-value sorting does not select the physical normal or inverted mass-eigenstate
labels.
\end{theorem}

\begin{theorem}[Exact principal-branch selector from the centered eta-class]
Fix the cycle sum
\[
\Omega=\psi_{12}+\psi_{23}+\psi_{31},
\]
and a positive weight family \(\mu_e\). Define
\[
\mu_{\min}=\min_e\mu_e,
\qquad
F_{\max}=\sum_e\arcsin\!\left(\frac{\mu_{\min}}{\mu_e}\right).
\]
If \(|\Omega|<F_{\max}\), then on the principal branch
\(\psi_e\in(-\pi/2,\pi/2)\), the affine energy
\[
A(\psi)=\sum_e \mu_e(1-\cos\psi_e)
\]
has a unique minimizer. It is given by
\[
\psi_e^\star=\arcsin(\lambda/\mu_e),
\]
where \(\lambda\) is the unique solution of
\[
\sum_e \arcsin(\lambda/\mu_e)=\Omega,
\qquad
\lambda\in(-\min_e\mu_e,\min_e\mu_e).
\]
\end{theorem}

\begin{proof}
The Euler--Lagrange equations are \(\mu_e\sin\psi_e=\lambda\). On the principal branch the scalar
function
\[
F(\lambda)=\sum_e\arcsin(\lambda/\mu_e)
\]
is strictly increasing because
\[
F'(\lambda)=\sum_e\frac{1}{\sqrt{\mu_e^2-\lambda^2}}>0.
\]
Its range on \((-\mu_{\min},\mu_{\min})\) is \((-F_{\max},F_{\max})\), so
\(F(\lambda)=\Omega\) has a unique solution under the stated condition. Strict
convexity follows because the Hessian of \(A\) is
\[
\nabla^2A=\mathrm{diag}(\mu_e\cos\psi_e),
\]
which is positive definite on the principal branch and is positive definite after
restriction to the affine plane \(\sum_e\psi_e=\Omega\).
\end{proof}


For the selected phases put
\[
 M(\eta)=\begin{pmatrix}
 a&\rho e^{i\psi_{12}}&\rho e^{i\psi_{31}}\\
 \rho e^{i\psi_{12}}&a&\rho e^{i\psi_{23}}\\
 \rho e^{i\psi_{31}}&\rho e^{i\psi_{23}}&a
 \end{pmatrix}.
\]

\begin{theorem}[Exact intrinsic spectral cubic]
Let \(H=M^\dagger M\). Then
\[
H=dI+T,
\qquad
d=a^2+2\rho^2,
\]
where \(T\) has zero diagonal and off-diagonals
\[
x_{12}=2a\rho\cos\psi_{12}+\rho^2e^{i(\psi_{23}-\psi_{31})},
\]
\[
x_{23}=2a\rho\cos\psi_{23}+\rho^2e^{i(\psi_{31}-\psi_{12})},
\]
\[
x_{13}=2a\rho\cos\psi_{31}+\rho^2e^{i(\psi_{23}-\psi_{12})}.
\]
Define
\[
P=|x_{12}|^2+|x_{23}|^2+|x_{13}|^2,
\qquad
Q=\Re\!\bigl(x_{12}x_{23}\overline{x_{13}}\bigr).
\]
Then the eigenvalues \(\lambda_k\) of \(T\) are exactly the three real roots of
\[
\lambda^3-P\lambda-2Q=0,
\]
and the intrinsic squared masses are
\[
m_k^2=d+\lambda_k.
\]
If \(P>0\), equivalently,
\[
\lambda_k=
2\sqrt{\frac{P}{3}}
\cos\!\left(
\frac13\arccos\!\left(\frac{3\sqrt3\,Q}{P^{3/2}}\right)-\frac{2\pi k}{3}
\right).
\]
If \(P=0\), all three off-diagonal entries vanish, \(T=0\), and
\(\lambda_0=\lambda_1=\lambda_2=0\); the trigonometric expression is not used.
\end{theorem}

\paragraph{Numerical mass readout.}
The cubic identity is an exact algebraic statement, not a stability claim
for its floating-point evaluation. Near a massless locus, subtracting
\(d+\lambda_k\) loses small positive squared masses; forming
\(M^\dagger M\) also squares the condition number. For example, at
\(\Omega=0\), \(\eta=0\), \(\rho=1\), and \(a=1+10^{-8}\), the two
small masses are \(a-\rho\), although binary64 evaluation of the cubic can
round both to zero. The executable mass readout uses a scaled direct
singular-value decomposition of \(M\), checks the positive Takagi congruence,
and retains the cubic calculation only as a diagnostic. These are normwise
floating-point residual checks, not interval enclosures or relative-error
certificates for arbitrarily small masses. A degenerate block requiring
additional congruence resolution is explicitly refused. This numerical
repair supplies no physical mass-eigenstate labels or source closure.

\begin{corollary}[Intrinsic eta-chain spectral closure]
Once the centered eta-class is emitted at the flavor boundary, the intrinsic neutrino branch
emits
\[
s_0\leq s_1\leq s_2,\qquad
s_1^2-s_0^2,\ s_2^2-s_0^2,\ s_2^2-s_1^2,
\]
and the corresponding intrinsic spectral subspaces with no PMNS import and no flavor-label
leakage. A columnwise Takagi basis up to signs is emitted only on the strictly
positive simple-spectrum locus; a simple zero mode retains its phase freedom. The physical normal-ordering assignment is
\((\nu_1,\nu_2,\nu_3)=(s_0,s_1,s_2)\), while the inverted-ordering assignment is
\((\nu_3,\nu_1,\nu_2)=(s_0,s_1,s_2)\). The source does not choose between them
and therefore does not select a physical ordering.
\end{corollary}

The scalar-certificate statement follows because centering determines all
\(\mu_e\), the principal-branch theorem determines the phases, and the
matrix and its spectral projectors then follow. The cubic is obtained by
expanding \(\det(\lambda I-T)\). These implications do not determine the
supplied reference, its absolute scale or the physical charged basis. On
the isotropic branch two singular values coincide, giving zero solar
splitting. In the anisotropic comparison the template-derived charged basis
gives angles incompatible with the adopted oscillation data
\cite{nufit61}; this is a failed conditional diagnostic rather than an OPH
prediction of PMNS angles. An ordering and a charged-sector intertwiner must
be selected before a physical neutrino comparison.

\section{Electroweak charts, scale and pole attachment}

\subsection{What the two-output chart can fit}

Fix positive running couplings \(\alpha_2,\alpha_Y\), a Higgs chart value
\(v\), and a source coefficient \(\eta\). In this subsection \(v\) uses the
conventional \(v_{\rm SM}\) normalization. Define
\[
 \tau_Y(\tau)=-\frac{\tau+2\eta}{1+4\tau^2},\qquad
 n_{\rm fib}(\tau)=1+\frac{\alpha_2\tau+\alpha_Y\tau_Y(\tau)}
                              {\alpha_2+\alpha_Y},
\]
\[
 M_W=v\sqrt{\pi\alpha_2(1+\tau)},\qquad
 M_Z=v\sqrt{\pi(\alpha_2+\alpha_Y)[n_{\rm fib}(\tau)+\delta n]}.
\]
Here \((\tau,\delta n)\) are chart coordinates, rather than emitted physical
self-energies. Any positive target pair in this chart can be fitted by
\[
 \tau=\frac{M_W^{\dagger2}}{\pi v^2\alpha_2}-1,\qquad
 \delta n=\frac{M_Z^{\dagger2}}{\pi v^2(\alpha_2+\alpha_Y)}-n_{\rm fib}(\tau).
\]
Substitution proves the exact hit. It is a measured-pair reparametrization.

\begin{proposition}[Local nondegeneracy of the electroweak chart]
On \(1+\tau>0\), \(n_{\rm fib}+\delta n>0\),
\[
 \det\frac{\partial(M_W,M_Z)}{\partial(\tau,\delta n)}
 =\frac{M_WM_Z}{4(1+\tau)(n_{\rm fib}+\delta n)}>0.
\]
\end{proposition}
\begin{proof}
The first mass is independent of \(\delta n\); the diagonal derivatives
are \(M_W/[2(1+\tau)]\) and \(M_Z/[2(n_{\rm fib}+\delta n)]\).
\end{proof}
Local invertibility does not prove that the source has exactly two response
coordinates or selects a particular point.

\begin{definition}[Finite quotient-path certificate]
\label{def:particle-d10-qt}
After removing gauge, port-label and scheduler presentation coordinates,
assume the CP-even charge-preserving response through two returns has
exactly two real scalar channels and no off-block term. Require independently
specified carrier paths, a primitive response
\(\lambda=\eta\alpha_U/4\), charged incidence \((1,2/3,1)\), neutral
incidence \((1,4/3,2)\), and a diagonal central projector of normalized trace
\(1/6\) and amplitude \(\rho\). Its degree-two subtraction is
\((\rho/6)\eta^2\) in both channels. Require contracting charged and
uplifting neutral signs, and a strictly convex hypercharge least-norm
problem with Gram factor \(1+4\tau^2\) and residual \(\tau+2\eta\).
The paths must exhaust the admitted class; deeper paths and alternative
weights must be absent, quotient-trivial or excluded by a positive,
target-independent selector gap.
\end{definition}

\begin{theorem}[Conditional quotient-transport value law]
\label{thm:particle-conditional-d10-qt}
Assume one target-independent source branch emits
\((\alpha_U,\alpha_2,\alpha_Y,\eta,v)\), with
\(\rho=(\alpha_2-\alpha_Y)/(\alpha_2+\alpha_Y)\), \(\eta=\rho\alpha_U\).
Under Definition~\ref{def:particle-d10-qt}, the unique admitted transport is
\[
 \tau=-\lambda[1+2\eta/3+(1-\rho/6)\eta^2],\qquad
 \delta n=\lambda[1+4\eta/3+(2-\rho/6)\eta^2],\quad
 \lambda=\eta\alpha_U/4.
\]
It gives the displayed chart masses whenever
\(\alpha_2'>0,\alpha_Y'>0\), where
\[
 \alpha_2'=\alpha_2(1+\tau),\qquad
 \alpha_Y'=\alpha_Y[1+\tau_Y(\tau)]+(\alpha_2+\alpha_Y)\delta n.
\]
\end{theorem}
\begin{proof}
The primitive normalization, return incidences, central subtraction and
signs fix the two polynomials. Strict convexity fixes the hypercharge
minimizer. Substitution gives the mass chart, and exhaustion excludes
alternative admitted responses.
\end{proof}
The certificate, including its exhaustive path category, is a hypothesis.
Assigning weighted paths to the monomials of a desired polynomial does not
prove that the physical source selects those paths or weights.

\subsection{Electromagnetic readout and Maxwell normalization}

\begin{theorem}[Conditional Ward-projected transport]
\label{thm:ward_u1q_transport}
Assume a common physical branch supplies the photon and neutral/charged
polarization kernels, a Ward projector onto \(Q=T_3+Y\), a positive
heat-kernel parameter \(t_Q(q^2)\), and a finite strictly positive Thomson limit.
If \(p_n\propto\exp(-t_Qn^2)\) and \(g_Q^2=t_Q/(2\pi)\), then
\[
 \alpha_{\rm em}^{-1}(q^2)=\frac{8\pi^2}{t_Q(q^2)},\qquad
 \Delta\alpha(q^2;m_Z^2)=\frac{t_Q(m_Z^2)}{t_Q(q^2)}-1.
\]
For \(a_0=\alpha_{\rm em}^{-1}(m_Z^2)\), the Thomson inverse is
\(a_0t_Q(m_Z^2)/t_Q(0)\).
\end{theorem}
\begin{proof}
Use \(\alpha=g_Q^2/(4\pi)\) on the same branch and take its stipulated
Thomson limit. The Ward projection separates the unbroken electromagnetic
line; it does not construct its kernel or the heat-kernel law.
\end{proof}

\begin{corollary}[Conditional Maxwell normalization]
\label{cor:ward_maxwell_normalization}
If this branch supplies a nonzero Maxwell kinetic action with coupling
\(g_Q\), varying it gives \(d{*F}=g_Q^2{*J}\), with \(dF=0\).
\end{corollary}
\begin{proof}
The field equation is the action's Euler--Lagrange equation; the Bianchi
identity follows from the curvature definition. The heat parameter alone
supplies no Maxwell action.
\end{proof}

\subsection{Hierarchy and common screen load}

The conditional transmutation branch uses
\[
 \frac{v}{E_{\rm cell}}=\exp[-\pi/(2\alpha_U(P))].
\]
Its weak-doublet multiplicity is \(N_c+1=4\), with additive load \(4P\).
Finite exterior multiplicity is insufficient to prove a physical equality
of screen and electroweak load carriers.

\begin{theorem}[Conditional screen/electroweak composition]
\label{thm:ew-tick-projection}\label{thm:screen-sieve-capacity-projection}
Assume twelve equal-trace port-center atoms, the integer total-twelve
counting grammar, normalized readback cost, oriented \((12,30,20)\)
icosahedral incidence and the stated refinement cocycle. An equal-weight
load \(X\) has local readout \(X/12\). For
\(X=\log(N/\pi)\), cell entropy \(P/4\) and multiplicity four, the defined
screen composition is
\[
 \Gamma_{\rm scr}=\frac P{12}\log(N/\pi).
\]
If an independently supplied correctable-record capacity \(N\) and physical
common-load map identify this with \(\log(E_{\rm cell}/v)\), then
\[
 \mathcal B_{\rm EW}(P,N):=\alpha_U(P)\log(N/\pi)-6\pi/P=0,
 \qquad N=\pi\exp[6\pi/(P\alpha_U(P))].
\]
\end{theorem}
\begin{proof}
The declared counting cost
\(H(q)=12+\sum_p(q_p-1)^2\) has its unique all-one minimum with gap two
on the integer total-twelve fiber. The oriented incidence supplies inverse
pairing, the proper \(A_5\) action and a rank-three frame with \(G^2=4G\)
\cite{ophmicrophysics}. Equal-weight additivity gives \(X/12\).
Multiplying the three stated factors gives \(\Gamma_{\rm scr}\). The
additional physical identification and transmutation relation give the
last equation.
\end{proof}
The integer Lyapunov cost, normalized readback cost and A3 Hessian are
different objects. This local lineage theorem does not construct the
federation nerve, identify a primitive observer or derive a global support
screen.

\begin{theorem}[Defined bridge-map fixed point]
\label{thm:ew-exact-capacity}
For fixed \(P\) and \(0<\lambda\le1\), the map
\[
 \mathcal C(P,x)=(1-\lambda)x+\lambda\frac{6\pi}{P\alpha_U(P)}
\]
is a contraction with unique fixed point
\(x_\star=6\pi/[P\alpha_U(P)]\). Therefore
\(\mathcal B_{\rm EW}(P,\pi e^{x_\star})=0\).
\end{theorem}
\begin{proof}
Its derivative is \(1-\lambda\); solve \(x=\mathcal C(P,x)\).
\end{proof}
This is a defined map, without a proof that it is the physical direct
record-capacity map. \label{rem:n-reserve-branches}
Finite-presence and Poisson reserve factors,
\(1-P/24\) and \(e^{-P/24}\), require different composition premises.
Neither follows from the contraction or from their shared local expansion.

\begin{proposition}[Declared-map RG/Higgs square]
\label{thm:rg-higgs-naturality}
If both routes through source and Higgs quotients are defined to preserve
the same tuple \((P,N,\Theta,F_{\rm crit}F_{\rm EW}(P,N))\), with
\(\Theta=(N/\pi)^{-P/12}\), their comparison defect is zero.
\end{proposition}
\begin{proof}
Both routes return the same tuple by definition.
\end{proof}
This proves no source selection of the maps or physical naturalness.
Finite completions sharing the stated antecedents can instead have defect
one. A commuting definition is not a calculation that physical relevant
deformations preserve the weak hierarchy.

\subsection{Operational scale and source Hamiltonian}
\label{wzh-source-closure-envelope}
\begin{proposition}[Conditional physical clock attachment]
If the same strict source branch emits the dimensionless energy
\(\varepsilon_{\rm clk}>0\) of a declared clock transition whose physical
frequency \(\nu_{\rm clk}\) defines the operational unit, then, with
\(h_{\rm P}=2\pi\hbar\) Planck's constant,
\[
E_\star=\frac{h_{\rm P}\nu_{\rm clk}}{\varepsilon_{\rm clk}}.
\]
Equivalently, \(\gamma_\star=\ell_\star\nu_{\rm clk}/c\) gives
\(E_\star=\hbar\nu_{\rm clk}/\gamma_\star\). An operator-norm perturbation
of at most \(\epsilon\) changes an isolated selected clock gap by at most
\(2\epsilon\). This is a conversion and stability theorem; the required
physical reference transition and SI binding are absent.
\end{proposition}

\begin{proposition}[Exact finite source-Hamiltonian gap]
\label{prop:finite-source-hamiltonian-gap}
Let \(D_\sigma\) be the sign-twisted seam derivative on the connected
observer-visible finite local domain, and let
\(H_\sigma=D_\sigma^\ast D_\sigma\) use the declared unit-counting measure.
The signed-incidence formula makes \(H_\sigma\) integer symmetric and positive
semidefinite. A kernel vector would supply a sign-consistent coloring on each
connected component. The declared reversing seam transport frustrates every
triangle, so the connected complex has no such coloring and the signed
incidence rank is \(8{,}662\). Therefore \(\ker H_\sigma=0\). With maximum
degree \(12\), the determinant bound gives
\[
  \lambda_{\min}(H_\sigma)\geq 24^{-8661}>0 .
\]
A sparse numerical refinement gives
\(\lambda_{\min}=0.11753670336118684\) with relative residual
\(1.49\times10^{-14}\). The decimal is measured rather than certified. The
exact statement is strict positivity and the displayed floor.
\end{proposition}

The determinant floor follows because the positive integer determinant
is at least one and every eigenvalue is at most 24. The numerical decimal
has a residual check, not an interval enclosure. This gap does not select a
physical field fiber, laboratory transition or determinant attachment.

For an isolated cesium ground manifold, set \(P=P_3+P_4\),
\(Q=1-P\), with \(\dim V_3=7\), \(\dim V_4=9\), and define
\[
 \mathcal F_P(z)=P(\widehat H-z)P-
 P\widehat HQ[Q(\widehat H-z)Q]^{-1}Q\widehat HP.
\]
\begin{proposition}[Cesium channel scalarization and monotone clock roots]
\(z\in\operatorname{spec}\widehat H\) exactly when
\(0\in\operatorname{spec}\mathcal F_P(z)\), with matching multiplicities.
Rotational invariance and the multiplicity-free decomposition
\(V_3\oplus V_4\) force \(\mathcal F_P(z)=f_3(z)P_3+f_4(z)P_4\), and
\(\mathrm d\mathcal F_P/\mathrm dz\preceq-P\) gives \(f_F'(z)\le-1\). One
sign bracket per channel therefore isolates exactly one level, and the clock
normal form is
\(H^{\mathrm{eff}}_{\mathrm{clock}}=E_0P+\widehat a_{\mathrm{Cs}}\,\mathbf I\cdot\mathbf J\)
with \(E_0=(9E_4+7E_3)/16\) and
\(\varepsilon_{\mathrm{Cs}}=E_4-E_3=4\widehat a_{\mathrm{Cs}}\).
\end{proposition}

The Schur complement preserves the isolated spectrum; Schur's lemma
scalarizes the rotational blocks. Differentiating the resolvent gives
\(\mathcal F'_P\preceq-P\), proving the stated root uniqueness after
sign brackets are supplied. No atomic Hamiltonian is derived by this
reduction. The electromagnetic convention, electron ratio, nuclear packet,
scalar evaluators, refinement tails and target-independent provenance are
needed to calculate the gap. Source extensions differing in an unfixed
parent can give different unique gaps.

\begin{proposition}[Operational scale metrology]
In SI units
\(E/\mathrm{GeV}=[E/(h\nu_{\mathrm{Cs}})]\,[h\nu_{\mathrm{Cs}}/\mathrm{GeV}]\),
and the second factor is an exact unit conversion. A source-only numerical
GeV mass is therefore equivalent to a source-only dimensionless operational
clock ratio. Replacing cesium by another clock species adds the burden of
predicting that clock's frequency ratio to the defining cesium transition.
\end{proposition}

A clock-gap coordinate chosen to reproduce a gravitational coupling is
a calibration checksum. Source-action maximum entropy gives a unique law
only after its physical operator basis and source moments are supplied;
that uniqueness principle does not supply the missing moments.

\subsection{Renormalization transport and physical singularities}
\begin{proposition}[Unique frozen RG and threshold transport]
Let the canonically normalized coefficient vector \(X(\mu)\) obey a locally
Lipschitz beta system on each interval of a frozen finite threshold list.
Assume that the maximal solution from each incoming matched value exists and
remains finite on the full closed interval, and that each matching output lies
in the domain of the next beta system. Also assume deterministic matching maps, fixed loop order, field content,
renormalization scheme, threshold locations, and decoupling order. Then one
source initial condition determines one low-energy coefficient vector. If
\(L_j\) bounds the beta-system Lipschitz constant and \(K_j\) bounds threshold
map \(j\), perturbations obey the corresponding product of
\(e^{L_j\Delta t_j}\) and \(K_j\), plus the declared truncation and matching
remainders.
\end{proposition}

\begin{proof}
Picard--Lindel\"of gives local existence and uniqueness; the stated
full-interval non-blowup and domain hypotheses extend the unique solution to
every declared threshold.
Gr\"onwall bounds perturbations on each interval; applying each deterministic
Lipschitz threshold map and inducting through the ordered list gives the
stated composition.
\end{proof}

\begin{theorem}[Conditional complex-pole theorem]
For \(B=W,Z,H\), let \(\Gamma_B^{\rm phys}(s,\xi)\) be the BRST-complete
physical inverse two-point block, including every declared mixing field after
the Ward/Slavnov--Taylor projection, analytically continued to a declared
Riemann sheet, and set
\[
D_B(s,\xi)=\det\Gamma_B^{\rm phys}(s,\xi).
\]
Assume a reference determinant has exactly one simple physical zero inside a
certified contour, no zero lies on the contour, and the full determinant obeys
\(|D_B-D_{B,0}|<|D_{B,0}|\) there. In the neutral sector also assume the
Ward-protected photon line has been separated from the massive eigenvalue.
Then the enclosed pole is unique and stable. With the exact convention
\[
s_B=\left(M_B-\frac{i}{2}\Gamma_B\right)^2,
\qquad M_B>0,\quad\Gamma_B\ge0,
\]
it determines one mass and width. If a Nielsen identity has the form
\(\partial_\xi D_B=C_B^{(\xi)}D_B\), the simple pole is gauge-parameter
independent~\cite{gambino1999,grassiKniehlSirlin2001}.
For \(W\) and \(Z\), interpretation as a physical resonance additionally requires
a dressed BRST-invariant current amplitude whose Laurent coefficient at that
pole is nonzero; for \(H\), the analogous scalar-amplitude coupling is a
separate condition. A simple zero of the inverse-propagator determinant alone
does not establish either coupling.
\end{theorem}

\begin{proof}
Rouch\'e's theorem preserves the number of enclosed zeros under the certified
perturbation. Simplicity gives the implicit-function stability bound and a
locally simple determinant zero. Differentiating \(D_B(s_B(\xi),\xi)=0\) and using the Nielsen
identity gives \(ds_B/d\xi=0\). An invertible analytic field redefinition
multiplies \(D_B\) by a nonzero analytic factor and therefore preserves the
pole set.
\end{proof}

The populated physical self-energy kernels, analytic-sheet and contour
certificates, widths, residues and scheme uncertainties are absent from
the two-output mass chart. A label such as ``pole mass'' in a schema does
not supply these objects.

\begin{theorem}[Source-selection and physical-pole composition]
\label{ew-selector-rigidity}
Suppose the admitted source maps are fully specified modulo gauge and
trivial reparametrization. Assume the physical readout is constant on that
class or a target-independent selector has one minimizer with a positive
gap. Require an acyclic dependency graph excluding measured masses and
calibrated proxies from every source ancestor, including the specification
of the maps. If one source root and physical clock emit the canonical
coefficients, the quotient-path and criticality split certificates hold,
and the RG and physical-pole hypotheses above are satisfied, then the
composition determines a unique enclosed pole triple \((s_W,s_Z,s_H)\).
\end{theorem}
\begin{proof}
The class/selector hypotheses identify one readout. Source separation
follows from its complete ancestry. The unique root and scale determine
one coefficient packet, the RG theorem transports it uniquely, and the
pole theorem supplies one stable physical singularity per sector.
\end{proof}

\label{ew-repair-tuple-selection}\label{ew-conditional-envelope}
\label{ew-residual-attribution}\label{ew-anchor-source-branch}
The displayed constructions do not satisfy this conjunction. Generic
symmetry and analyticity permit families such as
\(\tau=-c\eta^2\), \(\delta n=d(1-\rho)\eta^2\), and criticality shifts
\((\pi_y,\pi_\lambda)\mapsto(\pi_y+a\eta^N,\pi_\lambda+b\eta^N)\), with
different outputs. Even a fixed running mass allows the different kernels
\(s-m_R^2\) and \(s-m_R^2-\epsilon\). Thus a chart, analytic dependence and
channel counting cannot select a physical pole. Numerical residuals must
be attributed to separate source, truncation, matching, clock and
experimental uncertainties. Scale variation is not a source covariance,
and componentwise intervals alone do not specify a joint likelihood.

\section{Hadrons, records and physical interpretation}

\subsection{Hadronic backend}

A source-only hadron prediction requires a source QCD parameter map and
quotient ensemble, a vacuum/transfer construction, physical currents and
spectral observables, and controlled continuum, volume, light-quark and
sampling errors. A two-current spectral measure can feed running-alpha and
vacuum-polarization calculations; it cannot supply four-current light-by-light
or transition amplitudes by itself. Published hadronic payloads are separate
empirical inputs.

For a declared positive Euclidean spectral representation,
\[
 C_X(t)=\int e^{-Et}\,d\rho_X(E)
\]
relates correlators to energy levels. Extracting a stable-channel mass also
requires excited-state and finite-volume bounds, an operational scale and
a continuum extrapolation. The seeded \(2+1\) family and deterministic
sampling contract are execution specifications. The retained surrogate
uses proxy pion and nucleon masses in constructed correlators. Its
receipt-to-evaluation replay tests infrastructure, without supplying
production unquenched lattice QCD, Dirac solves, baryon contractions or
source-derived hadron masses.

\subsection{Observer records and particle identity}

The finite record theorem constructs a central public record algebra under
its compatibility hypotheses. Given a declared algebra/state, instrument
and two-wing representation, the measurement theorem gives Born
probabilities, the selected L\"uders conditioning, repeated-read stability
under support-preserving evolution, and the stated CHSH bound
\cite{ophconsensus,ophobservers}. Those results do not identify an arbitrary
internal memory state with a particle or prove global observer continuation.

The finite CAR construction describes quantum excitations of a specified
source action. No identification as topological defects is established.
A physical species identification needs its spin/statistics, gauge charges,
propagation, interaction channels and observable spectral response to agree
on one field theory. Persistent record tokens can instead be stitched
without an implementation-ID rule under the hypotheses in the appendix.
That theorem establishes continuation in affine event charts. It does not
supply a geodesic action, a mass shell or a physical particle detector.

\subsection{Antiparticles and asymmetry}\label{subsec:baryogenesis-source}

A physical chiral-field attachment includes conjugate antiparticle states.
A cosmological matter asymmetry additionally requires a physical CP-odd
current, its state and generator, a clock, thermal history and charge map.
For an antisymmetric CP-odd winding cocycle \(\omega\), its candidate rate
has form
\[
 \dot\Theta_R=\theta_0\sum_{q,q'}p(q,T)L_T(q,q')\omega(q,q').
\]
A CP-symmetric state and generator make terms cancel against their CP
images, giving zero. A settled quotient or an iteration count supplies
neither a nonzero physical current nor its rate. The finite anomaly
coefficient \(k_R=\sum_{\psi\,\rm LH}r_\psi2T_2(R_\psi)\) is conditional
on the physical charged-field and record-charge assignment. It is not a
baryogenesis prediction.

\section{Requirements for a source-derived spectrum}

A physical effective description is a claim about compatible observer
statistics and interventions. It need not be a fundamental field or
spacetime placed behind those observations. The construction must supply
one fixed dictionary for preparations, clock readings, complete detector
outputs and subsequent responses, with controlled error over its stated
domain. Agreement of eigenvalues or of a few two-point values alone is
weaker than that contract.

A complete construction must determine the charged chiral inverse
propagator \(\Gamma^{(2)}(p;H)\) from the same repair/memory process that
produces the observer records. Its Higgs derivative, kinetic forms,
background and clock must be derived together. The observer's detector
law must identify the physical coefficients or singularities with the
empirical mass packet withheld. A continuum refinement then needs the
claimed chirality and species content, controlled interactions and an
appropriate mass scheme. For confined and charged sectors this includes
confinement and infrared qualifications rather than an assumed isolated
free-particle pole.

The next calculations have the following dependency order:
\begin{enumerate}[leftmargin=2em]
\item Specify the carrier/code, state, actual repair and coherent operations,
scheduler, controls, access and complete records. Test the seventeen/eighteen
current and rank-fifteen obstructions in that exact domain, and identify
the spatial and internal actions on it. An additional control, bath or
transport compensation remains a premise until implemented and justified.
\item Extract kinetic derivatives, charged gaps and Higgs vertices jointly
from one continuing finite process at several momenta and one declared
clock convention. Preserve the shared quark kinetic form, complementary
modes, unsuccessful detector outcomes and continuation after measurement.
This finite test can precede a complete interacting continuum theorem.
\item Test a source mechanism that escapes an explicitly identified
normalization or affine-family hypothesis. Record-persistence, localization,
radiative and coherent phase constructions are possible classes to examine,
not an exhaustive list. A hierarchy calculation and a physical CP/mixing
calculation have distinct acceptance tests; a raw nonzero phase is
insufficient. Failure excludes the specified candidate or proved class,
not all mechanisms on the carrier.
\item Establish compatible locality, state/vacuum, chiral measure and anomaly
control, dynamical gauge interactions and controlled volume, spatial and
occupation limits for that regime. A finite regulator or free-field success
does not discharge this stage. Each additional threshold changes the
running and matching obligations.
\item Freeze the source, observation map, calibration packet and uncertainties
before the validation comparison. Retain a hash, date, predicted observables,
decision rule, all tested alternatives and prior target exposure. Compare
ratios within a stated scheme/scale, then sector strengths, mixing and the
operational absolute scale. Nine charged Yukawas and the four independent
CKM parameters provide a charged-spectrum test; they do not by themselves
validate the full interacting observer theory. Independent interaction,
neutrino and clock tests retain their own obligations.
\end{enumerate}

Registration after exposure to known masses makes a reproducible postdiction;
it does not make the architecture historically target-blind. Numerical
replay tolerances, matching/truncation errors and experimental uncertainties
must be distinguished. A failed comparison is reported for the frozen rule
and its hypotheses; changing the rule creates a new recorded candidate.

Several errors can make a numerical comparison appear stronger than it is.
A selected twelve-dimensional subspace need not exhaust the available
currents. State cooling need not equal Hamiltonian projection. Conditional
event samples need not describe the complete instrument. A stationary
coefficient average need not equal correlated dynamical evolution. A Gram
matrix must be whitened with the vertex. Orientation reversal need not be
physical time reversal. Common rescalings cannot repair wrong ratios.
Mixed-scale quark coordinates and running electroweak parameters cannot be
compared directly with physical poles. A normwise residual cannot certify
arbitrarily small singular values. Target-informed architecture or discrete
selection remains target-informed even when runtime inputs are segregated.
These distinctions delimit the claims of every calculation above.

The finite gauge/matter maps are useful structural constraints on this
problem. They permit calibrated Standard Model charged-matter structure
and exact tests of proposed repair coefficients, while leaving the
physical source dynamics undetermined. The tested uncalibrated laws fail;
their exclusions identify restricted classes rather than a contradiction
of every OPH construction. The representation results, calibrated controls
and physical-spectrum derivation therefore have separate conclusions.

\section*{Reproduction}

The source operators, invariant projectors and Higgs tensors are implemented
in \path{model/observer_particles/}. The complete memory, gate and detector
replay is \path{model/common_source/native_family_protocol.py}; the
chiral-candidate, hierarchy-diagnosis and full-repair-response suites contain
the associated finite controls and exclusions. The working equations,
acceptance conditions and evidence index are in
\path{docs/research/OPH_PARTICLE_MASS_RESPONSE.md} and
\path{code/particles/PARTICLE_PROVENANCE_LEDGER.md}. These executable
receipts distinguish exact algebra, floating-point diagnostics and empirical
comparisons. A hash check verifies custody; it does not replace execution
or extend a theorem's physical scope.

\appendix
\section{Conditional modes and threshold kinematics}

Declared Maxwell, flat-background pure Yang--Mills and two-derivative
Minkowski Einstein--Hilbert actions have respectively two transverse,
\(2\dim G\) perturbative transverse, and two transverse-traceless classical
modes. A quantum particle additionally needs a positive-energy physical
Hilbert space and observable positive-residue spectral response. In
particular, confined QCD is not assigned asymptotic free gluons by its
perturbative mode count.

The complete thirty-seam source maps to
\(D_6=\{z\in\mathbb Z^6:\sum_i z_i\equiv0\pmod2\}\), with a
three-dimensional response-metric completion \cite{ophmicrophysics}.
Proper-carrier naturality and a unique positive normalized minimizer give
equal weights on sixty directed moves. Its spatial symbol has a
two-dimensional momentum-orthogonal fiber at nonzero momentum. If that
symbol is declared an oscillator stiffness under a nonzero coordinate
dilation, \(A_T''+\Lambda_a(k)A_T=0\), with
\(\omega_a^2=\Lambda_a\), \(\omega_a(0)=0\). This supplies neither Gauss
law nor a physical electromagnetic action, position or clock.

Its positive cosine symbol has an auxiliary root satisfying
\(|\Omega(k)-\Omega(p)|\le|k-p|\) in the chosen carrier chart.
A physical photon identification and Lorentz-invariant positive-mass
leptons would exclude photon decay into a pair under additive
energy-momentum conservation. Other dispersion tests require their
separate lepton, interaction and detector hypotheses. For leading
dimension-six corrections, the head-on collinear pair threshold depends
at positron share \(x\in(0,1)\) on
\[
 \delta_{\gamma,2}-x^3\delta_{+,2}-(1-x)^3\delta_{-,2}.
\]
Equal sharing therefore reads one combination of three coefficients; its
coefficient fiber is two-dimensional. A photon-only bound assumes an
independent Lorentz-invariant-lepton branch.

In the universal subcase with common hard-particle coefficient
\(\delta=-d=-a^2/20\), and a leading Lorentz-invariant soft photon, put
\(u=x(1-x)\in(0,1/4]\). The required soft energy is
\[
 \epsilon_{\rm req}(k,u)=\frac{m_e^2}{4ku}+\frac{3dk^3}{4}u.
\]
The constrained minimum is at equal sharing when
\(3dk^4\le16m_e^2\); above that transition it is attained at
\(u_*=m_e/(\sqrt{3d}\,k^2)\), with
\(\epsilon_{\rm env}=(\sqrt3/2)m_e k\sqrt d\). Differentiation gives
the minimizer; AM--GM gives the same envelope and excludes solutions below
it. The executable calculation and the finite theorem verify the
open-share equality witness and endpoint regime. These are conditional
kinematic windows. They do not supply opacity, propagated flux, an
air-shower interpretation or a physical attachment of \(a\).

\section{Finite source constraints on flavor}
\label{quark-native-response-constraints}

The icosahedral source constrains the carrier and repair dynamics before a
mass readout is chosen. Let $N_{i+},N_{i-}$ be the record counts at the two
ports of antipodal pair $i$, and center their totals
$b_i=N_{i+}+N_{i-}$ to obtain
\[
 W_5=\{b\in\mathbb R^6:\textstyle\sum_i b_i=0\}.
\]
For unit vectors $n_i$ along the six axes, define
\[
 T(b)=\sum_{i=1}^6 b_i(n_i n_i^{\mathsf T}-I_3/3).
\]
The icosahedral inner products give
$T^*T=(4/5)P_5$, where $P_5=I_6-\boldsymbol1\boldsymbol1^{\mathsf T}/6$.
Thus $W_5$ is the geometric quadrupole carrier. Its $A_5$ character has
norm one, so it is irreducible; the invariant normalized linear functional
on its endomorphisms is $\operatorname{Tr}/5$.

\paragraph{A geometric response carrier and its heat channels.}
Supply one unoriented face axis. Its normalizer is $D_3$, and exact
restriction gives
\[
 W_5|_{D_3}=\mathbf1\oplus\mathbf2\oplus\mathbf2.
\]
The group average is a canonical rank-one projector $P_0$ with matrix
entries $\pm1/6$ in pair coordinates. Its normalized trace is $1/5$;
the complementary response has dimension four and trace weight $4/5$.
This realizes the five-dimensional relabeling carrier geometrically.
Selecting the face-axis background and identifying this quadrupole
response with physical flavor remain additional steps. In particular,
the physical geometric triplet restricts as
$\mathbf3|_{D_3}=\mathrm{sign}\oplus\mathbf2$, a different representation.

There is no need to insert an arbitrary regular submodule to represent
classical $D_3$ heat probabilities as channels on $W_5$. Conjugation gives
\[
 \operatorname{End}(W_5)|_{D_3}
 =5\mathbf1\oplus4\,\mathrm{sign}\oplus8\mathbf2.
\]
The Hilbert--Schmidt Gram matrix of its six group-conjugation
superoperators is $24I_6+J_6$, which is positive definite. Their mixture
coefficients are therefore identifiable. If each of the three
involutions has jump rate $c_2\geq0$ and each of the two three-cycles
has rate $c_3\geq0$, the nontrivial eigenvalues of the positive heat generator
are $6c_2$ and $3(c_2+c_3)$. At time $t\geq0$, put
\[
 r=e^{-3(c_2+c_3)t},\qquad s=e^{-6c_2t};\qquad r^2\leq s\leq1.
\]
The transposition-only relation $s=r^2$ requires $c_3=0$ for $t>0$.
The endomorphism heat trace is $5+16r+4s$; the regular-carrier trace
$1+4r+s$ is a different observable. The full native
30-seam port Laplacian acts on $W_5$ as $6I_5$.
A $D_3$ jump process is consequently a further dynamical hypothesis,
not the native seam process expressed in another basis.

\paragraph{A typed color coefficient.}
Normalize fundamental color generators by
$\operatorname{Tr}(T_aT_b)=\delta_{ab}/2$. Then
$\sum_aT_a^2=(4/3)I_3$, and the quadratic color observable
\[
 A=\Bigl(\sum_aT_a^2\Bigr)\otimes P_0
 \quad\hbox{on }\mathbb C^3\otimes W_5
\]
has
\[
 \frac{\operatorname{Tr}A}{15}=\frac4{15},\qquad
 \frac{\operatorname{Tr}A^2}{15}=\frac{16}{45},\qquad
 \frac{\operatorname{Tr}e^{wA}}{15}=\frac{4+e^{4w/3}}5.
\]
This supplies a source realization of the linear coefficient $4/15$.
The first four scalar cumulants are
$4/15,64/225,256/1125,1024/50625$; higher response terms are constrained
once this particular observable and trace response are chosen.
The observable is not a projector or a POVM effect: its nonzero
eigenvalue is $4/3$. Its identification as a mass response is a
physical hypothesis, and these cumulants do not derive a Gaussian
law on arbitrary composite modules.

The contraction order matters. A color-invariant sharp projector on
$\mathbb C^3\otimes M$ has rank divisible by three. Twirling a rank-one
component projector instead produces the unsharp invariant effect
$(I_3/3)\otimes|m\rangle\langle m|$. It preserves the trace but divides
the Hilbert--Schmidt square by three. Applying an isotropic Gaussian
quadratic response after this twirl changes its coefficient accordingly;
averaging nonlinear component responses afterwards need not do so.
Neither source symmetry nor dimension counting fixes that choice.
The exact character, channel, Casimir and moment identities are independently
reconstructed by
\path{code/particles/flavor/verify_source_w5_response_constraints.py}.

\paragraph{Exact constraints from strict-descent seam repairs.}
On a seam $(i,j)$, the strict-descent settlement rule transfers one record
from $i$ to $j$ only if $N_i-N_j\geq2$. It preserves $S=\sum_iN_i$ and decreases
$V=\sum_iN_i^2$ by $2(N_i-N_j-1)>0$. Every signed-integer fiber
$\{N:\sum_iN_i=S,\ V\leq B\}$ is finite and forward invariant.
With the additional assumption of state-independent uniform proposals
among the 30 seams, order the states by $V$. The transition matrix is
triangular with diagonal
\[
 K_{NN}=1-\frac{a(N)}{30},
\]
where $a(N)$ is the number of admissible seams. Its eigenvalues are exactly
these diagonal entries, and every stationary measure is supported on
settled states. The positive eigenvalues of $I-K$ have ratios at most
30. This is a discrete proposal-count statement; neither a physical
clock nor equality between these rates and particle masses is assumed.
The settlement rule differs from the native seam rule of the executable
observer construction, which sets the two endpoint loads to the floor and
ceiling of their mean and assigns an odd remainder by a fair coin. At load
difference one the native rule waits or swaps, so its balanced class keeps
moving; the triangular spectrum and the matching calculation below apply
to the settlement rule only \cite{ophmodel,ophrepairlaw}.

A complete finite calculation is possible in the invariant sector
\[
 N_i\in\{0,1,2\},\qquad \sum_{i=1}^{12}N_i=12.
\]
An active seam joins a port with count two to a port with count zero.
Repair sets both counts to one and deletes both endpoints from the active
bipartite graph; it creates no new active edge. Complete repair schedules
are therefore precisely orderings of maximal matchings of that graph. This proves
that the minimum settling cost is the smallest maximal-matching size,
and the longest accepted schedule has the maximum-matching size.
Uniform fresh seam proposals give a random-greedy-matching absorption
law, without a Gaussian or diffusion approximation.

Exhaustive enumeration gives 73,789 states, 1,275 rotational orbits,
and 303 settled states in nine orbits. Independent enumeration of
13,133 maximal matchings over orbit representatives verifies the settling calculation; the largest
minimum cost is four, and the longest accepted schedule has six moves.
The nine settled orbits leave an eight-dimensional simplex of
$A_5$-invariant stationary mixtures. Symmetry and settling alone do not
choose its weights.

The endpoint classification has a larger domain than that transient
sector. Settledness implies $|N_i-N_j|\leq d(i,j)$ for graph distance
$d$. Every icosahedral vertex has distance profile $(1,5,5,1)$, hence
$\sum_jd(i,j)=18$ and
\[
 \frac S{12}-\frac32\leq N_i\leq\frac S{12}+\frac32.
\]
If $S=12k$ with integer $k$, integrality forces
$N_i\in\{k-1,k,k+1\}$. Uniform shifts do not change the quadrupole.
Consequently the 303 settled states and their four spectra exhaust
\emph{all signed-integer endpoints} at such a total, without an initial
three-level or positivity assumption. Strict descent makes every maximal
accepted repair sequence finite; fresh uniform proposals reach an
endpoint almost surely, whereas an unfair proposal schedule can stall.
The settled quadrupole eigenvalue multisets are exactly
\[
 \{0,0,0\},\quad
 \{-2/\sqrt5,0,2/\sqrt5\},\quad
 \{-1-1/\sqrt5,\,1-1/\sqrt5,\,2/\sqrt5\},
\]
and the negative of the last multiset. Their multiplicities among
settled states are respectively $63,180,30,30$. Thus nonzero quadrupoles
satisfy
\[
 \operatorname{Tr}Q^2\in\{8/5,16/5\},\qquad
 \frac{(\operatorname{Tr}Q^3)^2}{(\operatorname{Tr}Q^2)^3}
 \in\{0,9/64\}.
\]
Under the additional affine log-mass readout $\log m_j=a+bq_j$, a
nondegenerate ordered spectrum must consequently obey
\[
 \frac{\log(m_2/m_1)}{\log(m_3/m_2)}
 \in\{1,R,R^{-1}\},\qquad R=\frac{3/\sqrt5-1}{2}.
\]
This constrains a specified pointwise readout at every total divisible
by twelve, but does not select the physical total, endpoint or vacuum.
It does not apply the same finite menu to an arbitrary averaged response
operator of a stationary mixture.

The unrestricted total also admits a finite classification. Quotient
settled counts by uniform integer shifts and set their minimum to zero.
The graph diameter bounds their range by three. Range-three extrema
must be antipodal; the identity $d(a,j)+d(-a,j)=3$ forces the unique
height profile $N_j=3-d(a,j)$. Its antipodal pair totals are constant,
so its quadrupole vanishes. There are twelve such profiles. Every
other representative has entries in $\{0,1,2\}$ and minimum zero.
Exhausting those assignments gives 6,065 settled representatives.
Thus \emph{all} signed-integer settled source states, at arbitrary total,
reduce to 6,077 classes modulo uniform shift and exactly 31 quadrupole
spectra. Their full characteristic polynomials and total residues modulo
twelve are recorded in the executable certificate.

There is also a sharp norm bound on every integer settled state.
The two seams between any two antipodal pairs form a perfect matching.
The settled seam bound therefore gives $|b_i-b_j|\leq2$ for pair totals.
The variance bound for six numbers in an interval of width two yields
$\sum_i(b_i-\bar b)^2\leq6$. Hence
\[
 \operatorname{Tr}Q^2\leq\frac{24}{5}.
\]
Equality is attained by setting all ports of three antipodal pairs to
one and all ports of the other three pairs to zero. The bound extends to real assignments satisfying the same unit
edge-Lipschitz condition; the integer repair threshold alone does not
impose that condition on arbitrary real counts.

For every nonzero endpoint quadrupole, the shape invariant
$\chi=(\operatorname{Tr}Q^3)^2/(\operatorname{Tr}Q^2)^3$ belongs to the
following exact set of eleven values:
\[
\begin{gathered}
 0,\quad \frac5{192},\quad \frac{45}{686},\quad
 \frac9{64},\quad \frac16,\\
 \frac{487\pm189\sqrt5}{6000},\quad
 \frac{809\pm216\sqrt5}{29478},\quad
 \frac{12509\pm2808\sqrt5}{146334}.
\end{gathered}
\]
The value $1/6$ means a repeated eigenvalue and cannot yield three
distinct masses under the affine log readout. For its ordered adjacent
log-gap ratio $g>0$, centering a triple with gaps $1,g$ gives
\[
 \chi(g)=\frac{(g-1)^2(g+2)^2(2g+1)^2}
 {24(g^2+g+1)^3}.
\]
This function is invariant under $g\mapsto1/g$ and strictly decreases
from $1/6$ to zero on $0<g\leq1$. The ten admissible values therefore
give exactly nineteen possible log-gap ratios: one at $g=1$ and nine
reciprocal pairs. This finite restriction requires neither a chosen
source total nor a noise approximation. It continues to require the
stated individual-state mass readout; additive log scale, overall
response strength, up/down attachment and physical mass scheme are
not selected by the stationary-shape theorem.

For the actual accepted pair-total increment, rejection or deletion of
an event can only decrease the sum of its outer products. Since each
edge of the six-pair complete graph has two native seam representatives,
uniform proposals give the statewise bound
\[
 0\preceq\operatorname{Cov}(\Delta b\mid N)
 \preceq\mathbb E[\Delta b\Delta b^{\mathsf T}\mid N]
 \preceq\frac25P_5.
\]
On quadrupoles the covariance bound is $(8/25)I_5$ in Frobenius
coordinates. At every settled state the accepted covariance vanishes.
A symmetric proposal alphabet, Gaussian noise or a physical Poisson
clock must not be inferred from this bound. The finite enumeration,
matching identities and exact moment checks are reproduced by
\path{code/particles/flavor/verify_native_repair_flavor_constraints.py}.

\paragraph{A conditional driven record state.}
A definite probability law for the relative record state follows from an
additional preparation protocol using the source's existing atomic events.
After each complete settlement, choose one of the twenty-four signed port
events $(i,s)$, $s=\pm1$, uniformly, append that event, and settle again using
fresh uniform seam proposals. The choice of this renewal protocol is a
hypothesis; it is not a consequence of the source grammar merely admitting
the events. No mass target, fitted orbit weight, or physical time scale enters.

If $N$ is settled, the state during repair after this event has the form
$N+s e_j$. A repair is possible exactly along a seam from $j$ to $k$ with
$s(N_j-N_k)=1$, and moves the single defect from $j$ to $k$. Thus a particle
walks strictly downhill in the original integer heights, and a hole strictly
uphill. The graph diameter is three, so at most three repairs occur.
Every accepted move has actual source--sink difference two. Rejected
proposals change waiting counts but not the uniform choice among admissible
neighbors.

Normalize heights by subtracting their minimum. On the resulting $6,077$
settled relative configurations the protocol defines a finite Markov kernel
with $66,192$ nonzero entries. Removing a record at a global maximum leaves
the state settled. Repeating reaches the zero representative, and reversing
these direct events gives positive-probability paths back. The kernel is
therefore irreducible. Every renewal flips the parity of the total residue,
and each direct edge has a two-cycle, so its period is exactly two. There is
a unique positive stationary distribution and Ces\`aro convergence; ordinary
stepwise convergence from one parity class is not claimed.

Proper rotations give $136$ orbits and a strongly lumped kernel with $1,000$
nonzero entries. The unique orbit weights are computed as exact rational
numbers and independently certified on every one of the $6,077$ states.
Their total-residue marginal is $1/12$ on each residue modulo twelve.
The law also respects sign reversal. For the centered count quadrupole,
\[
 Q(N)=\sum_iN_i\left(n_i n_i^{\mathsf T}-\tfrac13I\right),
 \qquad
 \mathbb E_\pi Q=0,
 \qquad
 \mathbb E_\pi[Q\otimes Q]
 =\frac{\mathbb E_\pi\operatorname{tr}Q^2}{5}\,I_5.
\]
The independently checked moments are
\[
 \mathbb E_\pi\operatorname{tr}Q^2\simeq1.4647638612,
 \qquad
 \mathbb E_\pi(\operatorname{tr}Q^2)^2\simeq2.8795832702.
\]
All odd centered-count moments vanish. In particular,
$\mathbb E_\pi\operatorname{tr}Q^3=0$. The radial fourth-moment difference
from a centered isotropic five-dimensional Gaussian is
\[
 \mathbb E_\pi(\operatorname{tr}Q^2)^2
 -\frac75\bigl(\mathbb E_\pi\operatorname{tr}Q^2\bigr)^2
 \simeq-0.1241631665\ne0.
\]
These are moments of the declared record law; they cannot be transferred to
an independently postulated physical response without its identification.

There is also an exact cost identity. Write
$W(N)=\sum_i(N_i-S/12)^2$. A uniform signed preparation increases $W$ by
$11/12$ on average. Each accepted one-defect repair decreases it by two.
Since $W$ is bounded on the relative state space, stationarity gives
\[
 \boxed{\mathbb E_\pi[\text{accepted repairs per preparation}]=\frac{11}{24}.}
\]
The stationary process is not reversible. For example, the transition from
$e_{11}$ to $e_{10}+e_{11}$ has probability $1/20$, its reverse has probability
$1/6$, and the exact stationary edge flows are unequal. This is a driven
stationary record law, not an identified thermal equilibrium or quantum
vacuum.

Only shift-invariant observables have this stationary law. The unquotiented
protected total performs an unbiased walk on $\mathbb Z$ and has no
normalizable stationary distribution. Subtracting a common height is a
readout projection, not a physical gauge identification of different total
records. Any Higgs or mass normalization depending on the mean count
requires that additional physical input. Likewise, averaging an
$A_5$-equivariant $\operatorname{End}(P_3)$ readout gives a scalar by
irreducibility, even for a nonlinear readout. The distribution of individual
sampled spectra is not the spectrum of that averaged operator. The renewal
law fixes the orbit weights conditionally; it supplies no
quark--Higgs map or mass scheme by itself.


\paragraph{The full exterior Higgs response and its kinetic normalization.}
Let the declared matter carrier be $C\oplus W=\mathbb C^3\oplus\mathbb C^2$,
with orthonormal exterior monomials, and write $H=h_0e_3+h_1e_4$.
The normalized top-form contractions
\[
 B_u(Q,H,u^c)=[Q\wedge H\wedge u^c]_5,\qquad
 B_d(Q,\overline H,d^c)=[(\iota_{\overline H}Q)\wedge d^c]_5
\]
give, in the ordered color and weak basis,
\[
 U(H)=I_C\otimes\binom{h_1}{-h_0},\qquad
 D(H)=I_C\otimes\binom{-\overline h_0}{-\overline h_1}.
\]
Here the sign convention for $d^c$ can be absorbed into its independently
supplied channel coefficient.  Exact exterior and independent creation-operator
calculations give
\[
 U^\dagger U=D^\dagger D=\|H\|^2I_C,\qquad U^\dagger D=0.
\]
Thus the full doublet, including its charged direction, has unit-magnitude
color Clebsches.  The channel coefficients $g_u,g_d$ are not determined by
normalizing these invariant tensors.  Color covariance forces the color
intertwiner to be a scalar identity; the traceless three- and five-dimensional
parts of the internal color current cannot be read as family masses.

On the conditional family attachment $F\otimes M$, restricting every supplied
operator to $I_F\otimes A$ gives $Y_u=g_uI_F$ and $Y_d=g_dI_F$.
This restriction is preserved by products, adjoints, inverses, positive
functional calculus, and elimination of heavy modes with a family-independent
block decomposition.  It is a precisely bounded direct-contraction model,
not a classification of every source interaction.  It produces no hierarchy.
After converting the all-left Weyl coefficients to the four-component Dirac
Yukawa convention, a local chiral action with positive, Higgs-independent
kinetic matrices has normalized vertex
\[
 Y_{q,\mathrm{can}}=
 Z_Q^{-1/2}Y_q Z_q^{-1/2}/\sqrt{Z_H}.
\]
Here $Z_q$ is the right-handed kinetic metric.  For a literal all-left
bilinear the left normalization factor instead carries a transpose;
the two coefficient conventions must not be mixed.
Consequently a proposed record matrix $R>0$ with
$Z_Q=Z_q=R$ and $Y_q=g_qR$ cancels from the vertex:
$Y_{q,\mathrm{can}}=g_qI_F/\sqrt{Z_H}$.
Equivalently, $\det(pR-hg_qR)=\det(R)(p-hg_q)^3$.
This is a conditional identity; no source selection of $Y_q=g_qR$ is assumed.
An interacting source completion must therefore emit a Higgs derivative
that is distinguishable from its kinetic metric, together with the common
state and normalization.  For a differentiable quadratic kernel with blocks
$A,B,C,D$, that derivative after eliminating an invertible heavy block is
\[
 \partial_H(A-BD^{-1}C)=
 A_H-B_HD^{-1}C-BD^{-1}C_H+BD^{-1}D_HD^{-1}C.
\]
The finite carrier construction does not supply this Higgs-dependent
source kernel or a physical mass prescription.

The independent finite replays are
\path{code/particles/flavor/verify_native_record_renewal.py}
and \path{code/particles/flavor/verify_exterior_higgs_response.py}.


\subsection{Stabilizers and entropy-selected quadrupoles}

Residual symmetry alone does not select the mass ratios. Realize \(W_5\) as
the traceless symmetric real \(3\times3\) matrices. Invariance under a
threefold or fivefold rotation about \(n\) forces
\[
 A=\alpha(nn^T-I/3),
\]
with spectrum \((2\alpha/3,-\alpha/3,-\alpha/3)\). Those fixed points have a
double eigenvalue. The twofold fixed locus is
\[
 A=\begin{pmatrix}u&v&0\\v&w&0\\0&0&-u-w\end{pmatrix},
\]
which is three-dimensional, hence two-dimensional after overall scaling,
and admits simple spectrum. It leaves exactly enough continuous freedom to
carry the two independent mass ratios without predicting them. The Klein
four-group of three perpendicular twofold axes fixes the diagonal traceless
forms, one projective parameter with generic simple spectrum, and every
subgroup of \(A_5\) other than the trivial group, \(C_2\), and \(V_4\)
contains a threefold or fivefold rotation, so no residual stabilizer pins an
isolated simple-spectrum orbit. The
source-side family problem is therefore selection by a specific
screen-derived \(A_5\)-invariant potential.

The leading cubic entropy truncation has a global classification. A \(W_5\) port vector consists of six values
\(a_i\), each repeated on an antipodal pair, with
\(\sum_i a_i=0\) and \(2\sum_i a_i^2=1\). At a constrained extremum of
\(S_3=2\sum_i a_i^3\), the Lagrange equation makes every \(a_i\) a root of
one quadratic. Hence there are at most two distinct values, and enumerating
their multiplicity \(m=1,\ldots,5\) gives
\[
 S_{3,m}=\frac{3-m}{\sqrt{3m(6-m)}}.
\]
The global maximum is \(2/\sqrt{15}\), uniquely up to permutation, with one
antipodal pair against the other five. It is the \(C_5\)-axis orbit, and its
quadrupole spectrum is exactly
\((-2/\sqrt{15},-2/\sqrt{15},4/\sqrt{15})\). Thus the leading cubic
truncation selects a double eigenvalue and fails the required simple-spectrum
shape. Interpreting that degeneracy as two equal physical charged masses
requires the separately unproved \(W_5\)-to-family attachment
(\path{code/particles/runs/leptons/charged_entropic_branch_no_go.json}).

At finite band amplitude, the global minimizers of one coefficient-free
candidate---the quartic truncation of the relative-entropy expansion on
\(W_5\)---can also be classified exactly. Write
\(b_i=1+r a_i=12q_i\). Exact stationarity, second-variation, and face
factorizations show that every global interior minimizer has two coordinate
values and that no omitted face wins. With
\[
 r_c=\frac{32\sqrt{15}-20\sqrt6}{27}=2.7757656378\ldots,
\]
the strict record domain \(b_i>0\) selects the degenerate multiplicity-one
\(C_5\) orbit for \(0<r<r_c\), both multiplicity-one and multiplicity-two
orbits at \(r=r_c\), and the simple-spectrum multiplicity-two orbit for
\(r_c<r<\sqrt{24}\). The latter has centered eigenvalues proportional to
\((-(\varphi^4-1),-1,\varphi^4)\) and sorted-gap ratio \(\varphi\). For
\(\sqrt{24}\le r<\sqrt{60}\), the strict problem has no attained minimum;
its infimum lies on the boundary.

The quartic truncation also admits a simple-spectrum boundary branch. On the closed probability
simplex the global minimizer above \(\sqrt{24}\) has four pair weights zero
and two weights
\(b_\pm=3\pm\tfrac12\sqrt{r^2-24}\). Its quadrupole spectrum is proportional
to \((-2,1-d,1+d)\), where \(d=\sqrt{(r^2-15)/5}\), so its upper-to-lower
sorted-gap ratio is
\[
 R(r)=\frac{2d}{3-d},\qquad
 r^2=15+\frac{45R^2}{(R+2)^2}.
\]
It reaches the compare-only central charged-lepton ratio
\(R=1.889009\ldots\) at the target-attached value
\(r=5.061326\ldots\). This is not a prediction: eight of twelve port
probabilities then vanish, the Taylor expansion about the uniform record is
least controlled there, no source emits \(r\), and the physical
quadrupole-to-log-mass attachment is not supplied. The strict full-support
quartic minima are conditionally excluded as the observed shape. The
closed boundary remains a minimum of that truncated objective
(\path{code/particles/runs/flavor/entropy_w5_shape_certificate.json}).

\paragraph{Full relative entropy at fixed pair norm.}
The untruncated uniform-reference objective gives a sharper answer for the
same six-pair model. Set
\[
 p_i=\frac{1+r a_i}{6},\qquad
 \sum_i p_i=1,\qquad
 \sum_i p_i^2=\frac16+\frac{r^2}{72},\qquad p_i\geq0,
\]
and minimize \(D(p\Vert\mathbf1/6)=\sum_i p_i\log(6p_i)\), with
\(0\log0=0\). Pairing the equal antipodal port probabilities leaves this
divergence unchanged. For every \(0<r\leq\sqrt{60}\), all global minima are
the six permutations of
\[
 p_{\rm high}=\frac{1+5r/\sqrt{60}}6,\qquad
 p_{\rm low}=\frac{1-r/\sqrt{60}}6
 \quad\hbox{(five entries)}.
\]
This is the standard maximum-Shannon-entropy branch at fixed index of
coincidence~\cite{harremoesTopsoe2001}; its implication for the icosahedral
quadrupole is the application here.

For completeness, compactness gives a minimum. At an interior extremum
with \(0<r<\sqrt{60}\), the two constraint gradients are independent, and
stationarity gives
\(\log p_i+1=\lambda+2\mu p_i\). Strict concavity of
\(\log x-2\mu x\) allows at most two positive roots \(x<y\).
If the larger root occurs twice, the variation exchanging those two
coordinates is tangent to both constraints, while its constrained Hessian
coefficient is negative:
\[
 \frac1y-2\mu
 =\frac1y-\frac{\log y-\log x}{y-x}<0.
\]
Thus an interior minimum has one high and five low entries, and the two
constraints give the displayed values.

No nonvertex boundary point minimizes the full objective. If its positive
coordinates include two unequal values, adding a new weight
\(\epsilon>0\) at a zero coordinate and adjusting those two values by
\(O(\epsilon)\) preserves both constraints, by the implicit function
theorem. The change in the objective is
\(\epsilon\log\epsilon+O(\epsilon)<0\). If instead the support consists
of \(k\geq2\) equal weights \(1/k\), replace those weights by
\((1-\epsilon)/k\), add the new weight \(\epsilon\), and add and subtract
\(\delta\) at two old coordinates, where
\[
 \delta^2=\frac\epsilon k-\frac{k+1}{2k}\epsilon^2.
\]
The perturbed weights stay positive for sufficiently small \(\epsilon\),
both constraints hold exactly, and the old-coordinate contribution changes
by \(O(\epsilon)\); the same negative leading term excludes this boundary
point. At \(r=\sqrt{60}\), purity is one and only the six vertices are
feasible. At \(r=0\), the distribution is uniform and does not select a
normalized direction \(a\).

For \(r>0\), every selected normalized port direction has pair values
\((5,-1,-1,-1,-1,-1)/\sqrt{60}\), up to permutation. With six unit
antipodal-axis representatives \(n_i\), the forward twelve-port map is
\(Q(a)=2\sum_i a_i(n_i n_i^{\mathsf T}-I/3)\); its normalization obeys
\(n_i^{\mathsf T}Q(a)n_i=8a_i/5\). Its spectrum is
\[
 \operatorname{spec}Q(a)
 =\left(-\frac2{\sqrt{15}},-\frac2{\sqrt{15}},\frac4{\sqrt{15}}\right).
\]
The six branches are the antipodal vertex axes; each has the order-ten
dihedral stabilizer \(D_5\) in the rotational icosahedral group
(the \(C_5\)-axis terminology names its fivefold axis).
The perturbation quadrupole \(Q(ra)=rQ(a)\) is zero at \(r=0\) and
has a repeated eigenvalue at every positive feasible amplitude. Consequently,
no amplitude rule \(r=r(P)\) rescues three distinct family masses through
an attachment that reads only these quadrupole eigenvalues and treats
equal eigenvalues equally. This excludes that uniform-reference,
fixed-norm full-entropy selector with that readout, including its closed
boundary; it does not exclude OPH, other references or constraints,
different actions, or different physical readouts. The quartic boundary
formula remains valid for its own objective; its two-positive-weight branch
does not minimize the full entropy, although both objectives share the pure
endpoint. Neither a physical mass attachment nor numerical quark
masses follow from this result
(\path{code/particles/runs/flavor/full_entropy_w5_selector.json}).

The finite theorem does not supply an effective action on \(W_5\) with a
unique simple-spectrum minimizing orbit. The invariant algebra has one
quadratic, two cubic, and two quartic coefficients, so the
\(W_5\)-restricted Hessian scalar \(h_5\) at the homogeneous state is a
necessary local discriminator. A physical charged shape additionally
requires the physical \(A_5\) family lift, normed determinant-line descent
with kinetic factors, an interacting charged kernel, an operational scale,
and a source packet fixed before comparison.

\section{Perturbative intrinsic neutrino response}
For the isotropic matrix put \(\varphi=\Omega/3\),
\(m_d^2=|a-\rho e^{i\varphi}|^2\), and
\(\Delta_{\mathrm{atm,iso}}=|a+2\rho e^{i\varphi}|^2-m_d^2\).
The following expansion uses the branch on which the collective state is
above the doublet.
Write
\[
\psi_e=\varphi+\delta_e,
\qquad
\delta_{12}+\delta_{23}+\delta_{31}=0.
\]
At first order in the centered eta-class,
\[
\delta_e=-\tan\varphi\,\eta_e+O(\eta^2).
\]
Define
\[
\sigma^2=\frac23\sum_e \delta_e^2.
\]
Assume \(2a\cos\varphi+\rho\neq0\) and that the collective state is separated from the
isotropic doublet. Then the two ascending doublet eigenvalues satisfy
\[
s_{0,1}^2=m_d^2\mp 2a\rho|\sin\varphi|\,\sigma+O(\delta^2),
\]
so
\[
s_1^2-s_0^2 = 4a\rho|\sin\varphi|\,\sigma+O(\delta^2)
=4a\rho\frac{\sin^2\varphi}{|\cos\varphi|}
\sqrt{\frac23\sum_e\eta_e^2}+O(\eta^2).
\]

Under the declared normal-ordering hypothesis, the ascending atmospheric gaps obey
\[
\Delta m_{31}^2 = \Delta_{\mathrm{atm,iso}} + \frac12\Delta m_{21}^2 + O(\delta^2),
\qquad
\Delta m_{32}^2 = \Delta_{\mathrm{atm,iso}} - \frac12\Delta m_{21}^2 + O(\delta^2),
\]
so the first-order invariant atmospheric object is the collective-to-doublet centroid gap
\[
\Delta_{\mathrm{cent}}=
s_2^2-\frac12(s_0^2+s_1^2)
=
\Delta_{\mathrm{atm,iso}}+O(\delta^2).
\]
Its quadratic shift is
\[
\delta\Delta_{\mathrm{cent}}
=
-
a\rho\,\sigma^2
\frac{a(4\cos^2\varphi-1)+6\rho\cos\varphi}
 {2a\cos\varphi+\rho}
+O(\delta^3).
\]

The projective largest-mass right-singular direction, in the phase gauge that
approaches the real democratic vector at the isotropic point, obeys the
first-order deformation law
\[
u_3=
u+
\kappa
\begin{pmatrix}
\delta_{23}\\
\delta_{31}\\
\delta_{12}
\end{pmatrix}
+O(\delta^2),
\qquad
u=\frac1{\sqrt3}(1,1,1)^T,
\qquad
\kappa=
\frac{\sqrt3\,(2a\sin\varphi+3i\rho)}
 {9(2a\cos\varphi+\rho)}.
\]
This equation fixes only the complex line. If \(v_3\) denotes its normalized
right-hand side, the canonical Takagi representative used by the executable
construction is
\[
u_3^{\rm Tak}=e^{-\frac{i}{2}\arg(v_3^T M_\nu v_3)}v_3,
\qquad
(u_3^{\rm Tak})^T M_\nu u_3^{\rm Tak}=s_3>0,
\]
up to the unavoidable column sign. In particular, at the isotropic point the
real democratic vector is generally outside positive Takagi gauge.
Equivalently,
\[
u_3=
u-
\tan\varphi\,\kappa
\begin{pmatrix}
\eta_{23}\\
\eta_{31}\\
\eta_{12}
\end{pmatrix}
+O(\eta^2).
\]

The expansion requires a separated collective state and
\(2a\cos\varphi+\rho\ne0\). It concerns the conditional matrix above.
A weighted-cycle comparison has additional unselected basis permutation,
holonomy orientation, CP sign and absolute attachment. Enumerating both
stored orientations and all row/column relabelings gives no
normal-ordering-consistent matrix passing both adopted correlated
oscillation gates. This excludes a convention-only repair of that matrix,
without classifying source-derived charged-basis placements.

\section{Affine event-record continuation}\label{sec:h3-record-worldlines}

The following conditional construction separates event supports, observer
frames, clocks and gauge transport. It is included to specify exactly what
record continuation establishes before any particle or geodesic attachment.

The compact cap-normal theorem supplies the Lorentz-natural frame hyperboloid. It identifies sky directions with
\(q(\Omega)=(1,\Omega)\), represents oriented round caps by
\(n_C=(\cot\alpha,\csc\alpha\,\mathbf c)\), and maps each cap to a geodesic half-space in
\(H^3\). The source-causal continuum certificate separately supplies affine
physical event charts only after exact two-way cone/order agreement,
count--volume, manifoldlikeness, topology, and refinement convergence. This section consumes event
supports, observer clocks, chart transitions, gauge transport, and worldline stitching. Frame
estimation, event location, overlap descent, and worldline stitching are separate certificates.

Fix a curvature radius \(R_H>0\) and write
\[
H^3_{R_H}=\{X\in\mathbb R^{1,3}:\langle X,X\rangle_L=-R_H^2,\ X^0>0\},
\qquad
\langle X,Y\rangle_L=-X^0Y^0+X^1Y^1+X^2Y^2+X^3Y^3.
\]
The hyperbolic distance used for conditioned frame comparison is
\[
d_H(X,Y)=R_H\,\mathrm{arcosh}\!\left(-\frac{\langle X,Y\rangle_L}{R_H^2}\right).
\]
An event-record atlas is a tuple
\[
\mathcal A_H=
\left(
R_H,\{U_i,\chi_i,u_i,J_i,\vartheta_i\},\{G_{ji},F_{ji}\},E,\nabla,
\mathsf{Prov}_{R_H},\mathsf{Prov}_u
\right).
\]
Here \(\chi_i\) are affine event charts, \(u_i\in H^3_{R_H}\) is the selected local observer-frame
anchor when such an anchor is claimed, \(J_i\) are event-chart domains for record support, and
\(\vartheta_i:J_i\to T\) are observer-clock maps into a common comparison-time line \(T\).
\(\mathsf{Prov}_{R_H}\) records whether \(R_H\) is a unit convention, an imported physical scale, or
an independently derived OPH scale. The unit-chart convention alone supplies
no independent OPH scale.
\(\mathsf{Prov}_u\) records which observer, tetrad, or clock object selected the anchors \(u_i\).
The Lorentz part \(G_{ji}\in\mathrm{SO}^{+}(1,3)\) transports frames and tetrads; the full
Poincar\'e transition transports affine event coordinates. The map
\(F_{ji}\) is the sector or gauge transport functor on the record payload. These two holonomies
are deliberately separate: atlas holonomy tests geometry, and gauge holonomy tests transported
sector data. When two charts describe the same transported observer frame, the atlas additionally
requires
\[
G_{ji}u_i=u_j.
\]

\begin{theorem}[H3 geometric naturality]\label{thm:h3-geometric-naturality}
For every \(A\in\mathrm{SO}^{+}(1,3)\), the map \(X\mapsto AX\) preserves \(H^3_{R_H}\),
\(d_H\), geodesics, exponential and logarithm maps, parallel transport, Hausdorff distances between
compact frame supports, and Fr\'echet means whenever the mean is unique.
\end{theorem}

\begin{proof}
The defining equation of \(H^3_{R_H}\) and the formula for \(d_H\) depend only on the Lorentz
inner product and the future sheet. Elements of \(\mathrm{SO}^{+}(1,3)\) preserve both. The
Levi-Civita connection, geodesic exponential, logarithm, and parallel transport are functorial for
isometries, so the listed constructions commute with \(A\). Hausdorff distance and unique
Fr\'echet means are metric constructions, hence are preserved by any isometry.
\end{proof}

\subsection{Cap-response frame balls and affine event supports}

The Standard Model gauge paper supplies the conditional inverse theorem that turns calibrated modular cap
responses into conditioned observer-frame supports. For every descended record token \(i\) and
clock slice \(t\), the frame layer consumes a certificate
\[
R_i(C,t,O)=\omega_{i,O}\!\left(\sigma_t^{C,O}(M_{C,0,O})\right),
\qquad
y_{i,j}=F_j(X_i(t))+e_{i,j},
\]
on a compact frame domain \(\Omega\subset H^3_{R_H}\), with quantitative observability
\[
\alpha\bar d(X,Y)\le |F(X)-F(Y)|_W\le L\bar d(X,Y),
\qquad |e_i|_W\le\sigma_i(t).
\]
Given an \(\varepsilon\)-net and residual tolerance \(\tau_i(t)\), the imported theorem gives the conservative enclosure
\[
S_i(t)\subseteq B_H(\widehat X_i(t),r_i(t)),
\qquad
r_i(t)=R_H\left[
\frac L\alpha\varepsilon+\frac2\alpha\sigma_i(t)+\frac1\alpha\tau_i(t)
\right].
\]
Not every point in the displayed ball is asserted to fit the responses.
A frame estimate \(\widehat X_i(t)\) is a unique finite output only when
\(\Delta_{\mathrm{frame},i}(t)>0\). If this gap fails, the output is a frame ambiguity set or
support ball. This ball cannot serve as an event position. Event stitching consumes an affine
event support \(A_i(t)\) produced by the physical source-causal continuum and
its ancestry/placement receipts. Distinct event
supports at a shared clock slice require a certified separation in the affine event metric. A
conditioned-frame separation may also be recorded:
\[
d_H(\widehat X_i,\widehat X_k)>r_i+r_k+m_{\mathrm{sep}}
\]
for a declared positive frame margin \(m_{\mathrm{sep}}\). It does not replace event separation.

\subsection{Proto-record extraction and descent}

A raw local proto-record is a connected component of an invariant detector scalar above a frozen
threshold, together with its certified affine event support, conditioned frame ball, local clock
interval, sector payload, and uncertainty collar.
The detector scalar is built from record-algebra data, not from implementation object IDs. Extraction is
chart-natural when transporting the raw field by \(G_{ji}\) and \(F_{ji}\) carries extracted
components in chart \(i\) to the extracted components in chart \(j\), up to the declared support
collar.

Before temporal stitching, overlapping charts are descended. The overlap-descent relation joins
two proto-records only when their transported affine event supports lie in the same connected
component on a real chart overlap and the same-component margin dominates the support collar. If distinct
component separation does not dominate the collar, the descent emits an ambiguity certificate
instead of choosing a label. On component-connected covers, this descent produces canonical
global descended records by the usual sheaf gluing argument: cocycle-compatible local components
with positive separation margins have a unique global component, and the construction is natural
under atlas restriction.

\subsection{Boundary-crossing germs and transported residuals}

A cross-boundary candidate edge is admissible only when the two descended records meet a real
patch interface in adjacent observer-clock intervals and share a common comparison chart. The
interface contact must be oriented and transverse: the signed-distance interval brackets the
interface, and the normal velocity has a positive lower bound. Tangencies, grazing contacts, and
file or shard-name boundaries without a physical interface are rejected. If the local slab contains
an interaction marker, the record is routed to the interaction solver instead of stitched by the
free-propagation rule.

For an admissible pair, the certificate transports the left record into the right comparison chart
and evaluates frozen residuals
\[
r_t,\quad r_{\partial},\quad r_x,\quad r_v,\quad r_{\mathrm{int}},
\]
for clock adjacency, boundary contact, affine event position, transported velocity or finite
difference, and interaction absence. The cost function is declared before seeing the proposed
matching. Record IDs, global IDs, shard-local IDs, and stitch keys are forbidden in admissibility,
costs, and tie-breaking.
All affine event-position and velocity residuals are evaluated as certified interval or
uncertainty-box distances. A point-estimate residual may be printed as a diagnostic; it cannot
replace the event-support radius in an admissibility or matching certificate. Conditioned
\(H^3\) frame residuals are a separate tetrad-consistency check.

\begin{lemma}[Candidate-edge naturality]\label{lem:h3-candidate-edge-naturality}
The admissible candidate-edge set and its residual intervals are invariant under common
\(\mathrm{SO}^{+}(1,3)\) chart changes and under gauge-bundle trivialization changes.
\end{lemma}

\begin{proof}
Every event term is computed from the affine overlap cocycle, oriented interface data, and the
common clock line. Frame terms use \(d_H\) and
Theorem~\ref{thm:h3-geometric-naturality}. Sector and gauge residuals compare transported payloads after applying the declared
connector, so changing a local trivialization conjugates both sides of the comparison. The
residual intervals and the resulting admissibility predicate are therefore unchanged.
\end{proof}

\subsection{Assignment gap, nonbranching paths, and ambiguity}

Let \(L\) and \(R\) be the descended record sets on two adjacent time slabs. A stitch matching is a
partial one-to-one assignment between \(L\) and \(R\), augmented by declared appearance and
disappearance interval costs. Its certified gap is
\[
\Delta_{\mathrm{cert}}
=\min_{M\ne M_*}\underline C(M)-\overline C(M_*),
\]
where \(M_*\) is the proposed matching, \(\overline C(M_*)\) is its upper cost bound, and
\(\underline C(M)\) is the lower cost bound for every competing matching.
The event-location gap \(\Delta_{\mathrm{event}}\) and stitch-assignment gap
\(\Delta_{\mathrm{cert}}\) are distinct receipts: the first certifies a unique event support
inside one clock slice, while the second certifies a unique temporal assignment across slices.

\begin{theorem}[Stable ID-independent stitch assignment]\label{thm:h3-stable-stitch-assignment}
If the candidate-edge set is natural, the costs are frozen and ID-independent, the proposed
matching is one-to-one, and \(\Delta_{\mathrm{cert}}>0\), then \(M_*\) is the unique certified
minimum-cost stitch. The certified matching is natural under chart changes, gauge
trivializations, and relabeling of implementation record IDs.
\end{theorem}

\begin{proof}
The strict gap says every competing matching has lower cost bound above the proposed matching's
upper cost bound, so no competitor can tie or beat it under any realization inside the certified
intervals. Lemma~\ref{lem:h3-candidate-edge-naturality} makes the candidate graph and cost
intervals invariant under chart and gauge presentation changes. Because IDs are excluded from
admissibility and costs, relabeling them changes only bookkeeping fields, not the optimum.
\end{proof}

\begin{theorem}[Unavoidable ambiguity]\label{thm:h3-unavoidable-ambiguity}
If a nontrivial permutation of the visible descended records preserves the atlas, clocks,
interfaces, sector/gauge payloads, candidate graph, and cost intervals, and its induced action
fixes no admissible stitch, then no natural ID-independent finite procedure can certify a unique
stitch across that slab.  Without the fixed-point-free clause, a unique natural stitch need only
be invariant under the symmetry and ambiguity does not follow.
\end{theorem}

\begin{proof}
A natural procedure must commute with the permutation. A unique certified output would therefore
be a fixed admissible stitch, contrary to the fixed-point-free hypothesis. Choosing one element
of a nontrivial orbit would introduce a label or presentation dependence. Thus
no natural single-valued admissible-stitch output depending only on the listed
preserved visible inputs exists. The declared fail-closed protocol records
this situation as \(\mathrm{AMBIGUOUS}\); an invariant orbit/set-valued output
or rejection is a different protocol and is not excluded.
\end{proof}

The certified stitch graph is then the graph obtained by chaining the unique matchings across
adjacent slabs. Since each matching is partial one-to-one, every connected component is a
nonbranching path, ray, or isolated record. These are the affine event-record worldlines of this
section.

\subsection{Refinement naturality and dynamical lift}

For a refinement \(s\to r\), require that record contraction induces a map
\(Q_{sr}:\mathcal M_s\to\mathcal M_r\) from every admissible fine matching
to an admissible coarse matching. The coarse optimum \(M_r^*\) must have
at least one admissible fine lift \(\widetilde M_s\), with
\(Q_{sr}(\widetilde M_s)=M_r^*\). Suppose the certified cost bounds imply
\[
|C_s(M)-C_r(Q_{sr}M)|\leq\eta_{sr}
\qquad(M\in\mathcal M_s)
\]
for all admitted realizations, and the coarse certified gap obeys
\(\Delta_r>2\eta_{sr}\). If a fine optimum \(M_s^*\) contracted to a
competitor, then
\[
C_s(M_s^*)\geq C_r(M_r^*)+\Delta_r-\eta_{sr}
>C_r(M_r^*)+\eta_{sr}\geq C_s(\widetilde M_s),
\]
a contradiction. Thus every fine optimum contracts to \(M_r^*\).
The lift premise is essential: coarse costs \(C_r(A)=0,C_r(B)=1\), a
single fine matching mapping to \(B\) with cost 1, and
\(\eta_{sr}=0\) satisfy the gap and distortion inequalities but cannot
recover \(A\). This is the refinement-naturality gate: feasibility,
the lift, and the uniform cost bound must all be certified before the
coarse/fine agreement follows from the margin.

The stitch certificate lives on the conditional effective Lorentzian manifold. A geodesic or
free-particle interpretation requires a separate dynamical receipt with a timelike speed margin,
connection compatibility, and an action or propagation law. The stitch theorem proves token
continuation in event charts; it does not prove geodesic motion.

\begin{theorem}[Certified affine event-record stitching]\label{thm:h3-record-worldline-stitching}
Given a source-selected refinement family satisfying the spacetime paper's
source-causal continuum certificate, together with a valid affine event atlas
compatible with its effective limit, a common
operational clock atlas, affine event supports for every descended record token and clock slice,
and separately typed conditioned-frame certificates where observer frames are used,
chart-natural proto-record extraction, overlap descent with positive event-support separation
margins, real transverse interface-crossing germs, sector and gauge transport continuity, a frozen
ID-independent one-to-one assignment with \(\Delta_{\mathrm{cert}}>0\), no interaction marker in
the free-propagation slab, and a passing coarse/fine contraction gate, the emitted stitch graph is
a unique natural nonbranching affine event-record graph. Under the declared fail-closed protocol,
failure to certify a positive event-location gap, separation margin, or assignment gap returns
\(\mathrm{AMBIGUOUS}\) and certifies no unique stitch by this gap contract. That protocol output is
not a proof of mathematical nonuniqueness; unavoidable ambiguity additionally requires the
fixed-point-free symmetry hypothesis of Theorem~\ref{thm:h3-unavoidable-ambiguity}. The protocol
returns \(\mathrm{INTERACTION\ REQUIRED}\) when an interaction marker is present, and rejects the
stitch when its atlas, interface, event-location, or transport gates fail.
\end{theorem}

\begin{proof}
The atlas, event-location, and extraction gates make local proto-records presentation-independent.
The positive \(\Delta_{\mathrm{event}}\) gate licenses event-position outputs; otherwise the
procedure carries event boxes or ambiguity sets forward. Descent turns overlap-compatible local components into
descended records before any temporal decision is made.
The interface and transport gates define a natural candidate-edge graph by
Lemma~\ref{lem:h3-candidate-edge-naturality}. The strict assignment gap gives a unique
ID-independent matching by Theorem~\ref{thm:h3-stable-stitch-assignment}. Chaining partial
one-to-one matchings gives nonbranching components, and the refinement gate makes the result
stable under the declared coarse/fine contraction. The other labels are the
declared protocol's fail-closed responses to failed or uncertified gates; they
need not be disjoint or logically exhaustive. Theorem~\ref{thm:h3-unavoidable-ambiguity}
upgrades the ambiguity label to a mathematical obstruction only under its
fixed-point-free action on admissible stitches.
\end{proof}

\section{Heat-Kernel Calibration Supplement}
\label{app:particle-d10-heat-kernel}

This appendix gives the compact-group heat-kernel calculation used by the
particle calibration. The fixed-cutoff same-overlap thermal/Casimir law is
proved in Ref.~\cite{ophmicrophysics}. The calculation below supplies the
compact-group Peter--Weyl lift and beta-shift bookkeeping. It does not derive
the physical one-loop beta coefficients from edge sectors alone.

\subsection{\texorpdfstring{\(\mathcal R_U\)}{R\_U} Unified-Coupling Certificate}

The \(\mathcal R_U\) certificate supplies the unified diffusion coupling used
by the hierarchy discussion of the flagship synthesis and by the electroweak
transmutation factor
\[
\frac{v}{E_\star}
=
P^{-1/2}\exp\!\left[-\frac{2\pi}{4\alpha_U(P)}\right].
\]
It does not certify the full clock branch
\[
\mathcal R_\gamma
=
\mathcal R_U+\mathcal R_\alpha+\mathcal R_e^{\mathrm{abs}}
+\mathcal R_{\mathrm{QCD/nuc}}^{133\mathrm{Cs}}
+\mathcal R_{\mathrm{atom}}^{133\mathrm{Cs}}.
\]
The complete \(\mathcal R_\gamma\) construction requires a source-only electromagnetic endpoint,
charged-lepton absolute scale, cesium nuclear data, and atomic spectral
enclosures. The numerical package is a witness for \(\mathcal R_U\). Its
mathematical conclusion requires the stated outward-rounded interval
arithmetic.

\paragraph{Declared source packet.}
The public-pixel branch uses
\[
P_C
=
\decnum{1.6309682094039593248792798477826489413359828516279250606661507533907793398933432}.
\]
The source-only branch uses the certified coordinate
\[
P_{\mathrm{fwd}}
=
1.630972095858897\ldots.
\]
It is interval-certified unique on its interval and, by the domain-global certificate, the only fixed point of its
readout map on the declared analytic domain \(\alpha^{-1}\in[100,200]\)
(the interval contraction certificate of 2026-07-14, whose closure row is a
certified source-root row).
The finite heat-kernel cutoffs are
\[
N_2=128,\qquad N_3=64.
\]
The one-loop running packet is
\[
(b_1,b_2,b_3)=\left(\frac{33}{5},1,-3\right),
\qquad
N_c=3,
\qquad
\beta_{\mathrm{EW}}=N_c+1=4.
\]
This packet is an explicit premise of the certificate. This appendix does not
derive it from the recovered OPH core.

\begin{definition}[No-hidden-free-variable rule for \(\mathcal R_U\)]
The following objects must be fixed before comparison with electroweak or gravity data:
\[
\begin{gathered}
(b_1,b_2,b_3),\quad
N_2,\quad
N_3,\quad
M_U(P),\quad
E_{\mathrm{cell}}(P),\\
\beta_{\mathrm{EW}},\quad
\hbox{Casimir normalization},\quad
\hbox{matching scheme}.
\end{gathered}
\]
Changing any of them after comparison with \(M_W\), \(M_Z\), \(\alpha_s(M_Z)\), low-energy
electroweak couplings, or \(G_{\mathrm{exp}}\) demotes the row to calibration.
\end{definition}

\paragraph{Dimensionless source map.}
Set \(E_\star=1\) for the dimensionless calculation. For a trial coupling \(a\), define
\[
M_U(P)=e^{-2\pi}P^{1/6},
\qquad
E_{\mathrm{cell}}(P)=P^{-1/2},
\]
\[
v(P,a)=E_{\mathrm{cell}}(P)\exp\!\left[-\frac{2\pi}{4a}\right].
\]
For \(i=1,2,3\),
\[
\alpha_i^{-1}(\mu;P,a)
=
a^{-1}
+
\frac{b_i}{2\pi}\log\!\left(\frac{M_U(P)}{\mu}\right),
\qquad
\alpha_Y(\mu;P,a)=\frac35\alpha_1(\mu;P,a).
\]
The source \(Z\)-scale satisfies the implicit equation
\[
\mu
=
\frac{v(P,a)}{2}
\sqrt{4\pi\alpha_2(\mu;P,a)+4\pi\alpha_Y(\mu;P,a)}.
\]
For \(\mathrm{SU}(2)\), with \(j=n/2\) and \(0\le n\le N_2\),
\[
d_j=2j+1,\qquad C_2(j)=j(j+1).
\]
For \(\mathrm{SU}(3)\), with \(0\le p,q\le N_3\),
\[
d_{p,q}=\frac{(p+1)(q+1)(p+q+2)}{2},
\qquad
C_2(p,q)=\frac{p^2+q^2+pq+3p+3q}{3}.
\]
The group sums are
\[
Z_G(t)=\sum_R d_R e^{-tC_2(R)},
\qquad
\bar\ell_G(t)=
\frac{1}{Z_G(t)}
\sum_R d_R e^{-tC_2(R)}\log d_R.
\]
Define
\[
t_2(P,a)=4\pi^2\alpha_2(\mu_Z(P,a);P,a),
\qquad
t_3(P,a)=4\pi^2\alpha_3(\mu_Z(P,a);P,a),
\]
and
\[
\Phi_U(P,a)
=
\bar\ell_{\mathrm{SU}(2)}(t_2(P,a))
+
\bar\ell_{\mathrm{SU}(3)}(t_3(P,a))
-
\frac{P}{4}.
\]
The certificate value is the zero \(\Phi_U(P,\alpha_U(P))=0\).

\paragraph{Numerical witness.}
For \(P=P_C\), the checker emits
\[
\begin{aligned}
\alpha_U(P_C)&=0.041124336195630495,\\
\alpha_U(P_C)^{-1}&=24.316501918546496.
\end{aligned}
\]
The corresponding dimensionless hierarchy data are
\[
\begin{aligned}
\frac{M_U(P_C)}{E_\star}&=0.0020260720012100037,\\
\frac{v(P_C)}{E_\star}&=2.0199803239725553\times10^{-17},\\
\frac{\mu_Z(P_C)}{E_\star}&=7.502403238986022\times10^{-18}.
\end{aligned}
\]
The emitted couplings at \(\mu_Z\) are
\[
\begin{aligned}
\alpha_1&=0.016885708038988166,\\
\alpha_Y&=0.010131424823392899,\\
\alpha_2&=0.033777889448908055,\\
\alpha_3&=0.11833602747445048.
\end{aligned}
\]
The heat-kernel parameters and entropy readouts are
\[
\begin{aligned}
t_2&=1.3334976254578108,\\
t_3&=4.6717191102770705,
\end{aligned}
\]
\[
\begin{aligned}
\bar\ell_{\mathrm{SU}(2)}&=0.39488144681077636,\\
\bar\ell_{\mathrm{SU}(3)}&=0.012860605540213493.
\end{aligned}
\]
Thus
\[
\begin{aligned}
\bar\ell_{\mathrm{SU}(2)}+\bar\ell_{\mathrm{SU}(3)}
&=0.40774205235098987,\\
P_C/4&=0.4077420523509898,
\end{aligned}
\]
and
\[
\begin{aligned}
\Phi_U(P_C,0.041124336195630495)
&=5.55\times10^{-17}.
\end{aligned}
\]
A centered derivative check gives
\[
\partial_a\Phi_U(P_C,a)\simeq -10.990524683118785.
\]
With \(\delta=10^{-6}\), the bracket signs are
\[
\Phi_U(P_C,\alpha_U-\delta)=1.0990748697592423\times10^{-5},
\]
\[
\Phi_U(P_C,\alpha_U+\delta)=-1.0990300680135956\times10^{-5}.
\]

For the certified source-only pixel,
\[
\begin{aligned}
P_{\mathrm{fwd}}&=1.630972095858897\ldots,\\
\alpha_U(P_{\mathrm{fwd}})&=0.041124247441816685\ldots,\\
\frac{v(P_{\mathrm{fwd}})}{E_\star}&=2.0198114078576331\times10^{-17}.
\end{aligned}
\]
This branch is the one to cite when avoiding the measured Thomson endpoint in the upstream pixel
solve.

\begin{theorem}[Electroweak unified-coupling numerical witness]\label{thm:particle-ru-witness}
Assume the declared source map, running packet, heat-kernel cutoff policy, and
no-hidden-free-variable rule above. Let
\[
I_U=[0.041123336195630494,\;0.041125336195630496].
\]
If outward-rounded interval arithmetic verifies the directed endpoint signs
\[
\Phi_U(P_C,0.041123336195630494)>0,
\qquad
\Phi_U(P_C,0.041125336195630496)<0,
\]
and verifies
\[
0\notin \partial_a\Phi_U(P_C,I_U),
\]
then there is a unique \(\alpha_U(P_C)\in I_U\) satisfying
\[
\Phi_U(P_C,\alpha_U(P_C))=0.
\]
The comparison pixel \(P_C\) inherits its CODATA electromagnetic input.
Holding that pixel fixed, this calculation uses no additional measured
low-energy coupling, measured \(G\), Planck unit formed using measured
\(G\), measured \(\Lambda\), measured \(M_Z\), or measured \(M_W\).
\end{theorem}

\begin{proof}
Continuity follows because the running couplings, self-consistent \(Z\)-scale solution, and
finite heat-kernel sums are continuous on the declared positive interval. The endpoint sign
change gives existence by the intermediate value theorem. The derivative exclusion gives
uniqueness. The dependency statement follows by inspection of the declared source map: the residual
is built only from \(P_C\), the declared running packet, the self-consistent \(Z\)-scale equation,
and finite compact-group representation sums.
\end{proof}

\begin{theorem}[Implicit-scale attachment at the comparison pixel]
\label{thm:particle-implicit-scale-attachment}
Let $F(a,\mu)$ be the right-hand side of the source $Z$-scale equation,
with the exact stored decimal
\[
 \widehat P_C=1.630968209403959324879279847782648941,
 \qquad J=[7\times10^{-18},8\times10^{-18}].
\]
For every $a\in I_U$, all running denominators are positive on $J$ and
\[
 F(a,J)\subset[7.49,7.52]\times10^{-18}\subset\operatorname{int}J,
 \qquad 0\le\partial_\mu F(a,J)\le0.006<1/100.
\]
There is exactly one scale $\mu(a)\in J$, it is continuously differentiable,
and the finite-cutoff hierarchy residual evaluated at this exact scale obeys
\[
 \frac{\mathrm d}{\mathrm da}\Phi_U(\widehat P_C,a)\in[-12,-10].
\]
The lower and upper endpoint values lie respectively in
\[
 [1.09907,1.09908]\times10^{-5},\qquad
 [-1.09904,-1.09902]\times10^{-5}.
\]
There is therefore exactly one zero in $I_U$. The cutoffs remain $N_2=128,N_3=64$; neither an
infinite representation sum nor a scale outside $J$ is part of this claim.
\end{theorem}

\begin{proof}
The uniform self-map and derivative bound give a contraction on $J$.
The interior image and $1-\partial_\mu F\ge99/100$ give the implicit derivative
\[
 \mu'(a)=\frac{\partial_a F}{1-\partial_\mu F}.
\]
Every interval image of a box containing this fixed point retains it.
Eight images on $I_U\times J$ enclose the scale for the derivative test;
sixteen images at each endpoint enclose the endpoint residuals. These are
inclusions of the exact fixed point, not assertions that an iterate equals it.
Interval automatic differentiation and an independent analytic calculation
verify the displayed bounds at 60-digit interval precision. For the latter,
write $S=\alpha_2+3\alpha_1/5$ and
$B=\alpha_2^2+99\alpha_1^2/25$. Then
\[
 F_\mu=\frac{FB}{4\pi\mu S},\qquad
 F_a=F\left[\frac{\pi}{2a^2}
 +\frac{\alpha_2^2+3\alpha_1^2/5}{2a^2S}\right].
\]
The finite entropy derivative is
$\mathrm d\bar\ell_G/\mathrm dt=-\operatorname{Cov}_t(C_2,\log d)$;
combining it with
$t_i'=4\pi^2\alpha_i^2[a^{-2}+b_i\mu'/(2\pi\mu)]$
gives the certified residual derivative. Continuity, endpoint signs, and
strict monotonicity finish the argument. The separate 60- and 64-iteration
certificates retain their own finite-map scope and Krawczyk enclosures.
\end{proof}

\paragraph{Precision policy.}
The exponential map makes downstream gravity displays sensitive to \(\alpha_U\):
\[
\frac{\partial\log G}{\partial \alpha_U}
\simeq
\frac{\pi}{\alpha_U^2}
\approx
1.86\times10^3.
\]
Therefore the number of digits printed for any downstream \(G\) or \(\varepsilon_{\mathrm{Cs}}\)
row may not exceed the interval precision certified by
\[
\mathcal R_U+\mathcal R_\alpha+\mathcal R_e^{\mathrm{abs}}
+\mathcal R_{\mathrm{QCD/nuc}}^{133\mathrm{Cs}}
+\mathcal R_{\mathrm{atom}}^{133\mathrm{Cs}}.
\]
The witness supports an electroweak hierarchy certificate. It does not support a 52-digit
source-only gravity prediction.

\subsection{One-Loop Calibration Frame}

At one loop, if couplings unify at \((M_U,\alpha_U)\),
\[
\alpha_i^{-1}(M_Z)=\alpha_U^{-1}+\frac{b_i}{2\pi}\ln\frac{M_U}{M_Z}.
\]
On the forward calibration branch these are the values
\((M_U(P),\alpha_U(P))\) emitted by the source solve. They are not inferred by
inverse readback from measured low-energy couplings.
Writing \(A_i:=\alpha_i^{-1}(M_Z)\) and \(L:=\ln(M_U/M_Z)\), one finds
\[
L=\frac{2\pi}{b_1-b_2}(A_1-A_2),
\]
and the corresponding consistency relation for this calibration is
\[
A_3^{\mathrm{pred}}
=
\frac{b_3-b_2}{b_1-b_2}A_1
+
\frac{b_1-b_3}{b_1-b_2}A_2.
\]
This appendix keeps that algebra on the page as a calibration relation, not as a theorem that OPH
derives a supersymmetric UV spectrum.

\subsection{Benchmark coefficient comparison}

For the standard one-loop benchmark coefficients,
\[
\begin{array}{c|c|c}
\text{Model} & (b_1,b_2,b_3) & \alpha_s(M_Z)^{\mathrm{pred}} \\
\hline
\text{Standard Model} & (41/10,-19/6,-7) & \approx 0.071 \\
\text{minimal supersymmetric Standard Model} & (33/5,1,-3) & \approx 0.116 \\
\text{Observed \cite{pdg2026}} & & \approx 0.118
\end{array}
\]
At one loop, the printed minimal supersymmetric Standard Model (MSSM)
coefficients reproduce the observed closure far better than the plain
Standard Model coefficients. The forward calibration matches this benchmark.
It does not imply that OPH derives a supersymmetric ultraviolet spectrum.
The observed row is the Particle Data Group world average, quoted at the
precision this benchmark separates, and the one-loop comparison is the
textbook unification benchmark rather than an output of the OPH source, which
is why it carries no claim class and no comparison row of its own.

\subsection{Heat-Kernel Sector Law}

Read this subsection as the compact-group merge boundary above the fixed-cutoff microphysics
theorem. MaxEnt yields the familiar \(d_R e^{-tC_2(R)}\) law only after the
quadratic-Casimir generator and one-sided edge algebra below are declared;
the Peter--Weyl lift alone does not select either input. This is not a second,
independent proof of the finite-cutoff same-overlap Casimir branch.

\begin{theorem}[Heat-kernel edge-sector weights]
Let the MaxEnt edge generator be the factorwise group Laplacian
\(H_{\rm edge}=\sum_i c_i\Delta_{G_i}\), with eigenvalue
\(C_2(R)\) in the one-factor notation, and let the probability space be the
one-sided electric-center algebra \(\bigoplus_R B(V_R)\), where the sector
projector has trace \(d_R\). Then the sector probabilities take heat-kernel
form:
\[
p_R(t)\propto d_R\,e^{-tC_2(R)},
\]
where \(d_R=\dim R\) and \(C_2(R)\) is the quadratic Casimir.
Bi-invariance by itself permits other central energies, such as
\(C_2(R)^2\). A Gibbs trace on the two-sided Peter--Weyl space \(L^2(G)\)
instead gives the multiplicity \(d_R^2\).
\end{theorem}

\begin{lemma}[Bi-invariant operators]
If \(G=\prod_i G_i\) is a compact semisimple Lie group written as a product of compact simple
factors, then any bi-invariant second-order differential operator on \(G\) has the form
\[
D=c_0\mathbf 1-\sum_i c_i\Delta_{G_i},
\]
where \(\Delta_{G_i}\) is the Laplace--Beltrami operator on the factor \(G_i\).
\end{lemma}

\begin{remark}
Bi-invariance is equivalent to \(D\in Z(U(\mathfrak g))\). Linear terms vanish because there are
no invariant vectors in the adjoint representation, and the quadratic part is proportional to the
Killing form on each simple factor, giving one independent coefficient per factor. The Casimir
element then acts as \(-\Delta_G\) in the regular representation.
\end{remark}

\subsection{Peter--Weyl Multiplicity and Beta Shifts}

By Peter--Weyl decomposition, the effective refinement-limit edge representation space is
\[
L^2(G)\cong \bigoplus_R V_R\otimes V_R^*.
\]
Entanglement traces over one factor give multiplicity \(d_R\) in \(p_R\), but loops see both
factors, so the effective multiplicity entering the calibration branch is
\[
N_{\mathrm{eff}}(R)=d_R\,p_R.
\]
The one-loop beta shift from edge sectors is then
\[
\Delta b_a=\sum_R p_R\,d_R\,T_a(R),
\]
where \(T_a(R)\) is the Dynkin index for gauge factor \(a\). This is the exact mathematical point
where edge-sector heat-kernel weights feed the calibration branch. Matching to
MSSM-like one-loop running is a comparison benchmark, not a theorem-level
MSSM spectrum claim.

The branch boundary for theorem-level unification closure is explicit. One requires a derived
refinement-limit identification of which transportable sectors actually contribute to the running
carrier, a derived statistics/grading rule rather than an imposed fermionic-grading restriction,
controlled threshold and decoupling conventions on that same carrier, a derived representation
content/truncation rule for the sectors retained in the comparison, and a proof that the effective
loop multiplicities entering \(\Delta b_a\) are exactly the ones used in this calibration package
rather than nearby alternatives. The declared runtime surface does not emit the
full RG/matching/threshold/scheme packet: scheme lock, threshold map,
beta provenance, and interval composition are separate certificate requirements.

At the branch unification diffusion parameter \(t_U(P)\approx 1.64\) selected
by the forward source solve, the benchmark shift is
\[
\Delta b\approx (2.49,4.38,3.97)
\qquad
\text{vs MSSM}
\qquad
(2.50,4.17,4.00).
\]
With the additional calibration assumption of a fermionic-grading restriction to
half-integer \(\mathrm{SU}(2)\) sectors, this sharpens to
\[
\Delta b\approx (2.50,4.17,3.97),
\]
which matches the MSSM-style one-loop benchmark at the percent level under that declared
calibration package. Equivalently, the benchmark ratio
\[
\frac{\Delta b_3}{\Delta b_2}\approx 0.91
\]
sits close to the MSSM comparison value \(0.96\).

\section{H3 Worldline-Stitch Certificate}\label{app:h3-worldline-certificate}

This appendix records the finite evidence-facing certificate for
Theorem~\ref{thm:h3-record-worldline-stitching}. It is separate from the heat-kernel
calibration package above. A passing certificate is a record-continuation statement in a declared
\(H^3\) observer chart, not a particle-species, mass, charge, or scattering-amplitude theorem.

The certificate payload has the following required blocks.
\begin{enumerate}
\item \textbf{Atlas.} It declares the hyperboloid model \(H^3_{R_H}\subset\mathbb R^{1,3}\),
the Lorentz signature \((-+++)\), the curvature radius \(R_H\), chart domains, and chart
transitions \(G_{ji}\in\mathrm{SO}^{+}(1,3)\). Transition residuals must certify Lorentz
inner-product preservation, orientation, and future-sheet preservation.
\item \textbf{Clock map.} Each local record interval is mapped to a common comparison-time line.
The observer-time adjacency margin must dominate the declared clock uncertainty.
\item \textbf{Chart-natural extraction.} Raw local records are connected components of a frozen
detector scalar on the \(H^3\) chart. The detector scalar, thresholds, support collars, and
chart-naturality residuals are included. Implementation record IDs are not inputs.
\item \textbf{Overlap descent.} Records are first descended across genuine chart overlaps. The
payload records same-component join margins, distinct-component separation margins, support
errors, and triple-overlap cocycle residuals.
\item \textbf{Interface crossing.} A candidate stitch must cross a real interface in adjacent
clock intervals. It includes oriented signed-distance margins and a positive lower bound on
normal velocity. Tangential, grazing, and file-boundary-only contacts are rejected.
\item \textbf{Transport.} The candidate edge is evaluated in a common chart after sector and
gauge transport. The atlas holonomy comparison is separate from the gauge-holonomy comparison.
\item \textbf{Assignment.} The matching is one-to-one with appearance/disappearance interval
costs. The proposed winner has an upper cost bound, every competitor has a lower cost bound, and
\[
\Delta_{\mathrm{cert}}
=\min_{M\ne M_*}\underline C(M)-\overline C(M_*)>0 .
\]
Record IDs, global IDs, stitch keys, and shard-local IDs are forbidden in admissibility, costs,
and tie-breaking.
\item \textbf{Refinement.} A coarse/fine pair supplies \(Q_{sr}\), a distortion bound
\(\eta_{sr}\), and a contracted-graph isomorphism check. The coarse gap must satisfy
\(\Delta_r>2\eta_{sr}\).
\item \textbf{Interaction firewall.} A free-propagation slab is required. If an interaction marker
is present, the terminal label is \(\text{H3 interaction required}\) and the event is handed
to the interaction solver.
\end{enumerate}

The allowed terminal outcomes are
\[
\begin{gathered}
\text{H3 stitch certified},\quad
\text{H3 stitch ambiguous},\quad
\text{H3 stitch rejected},\\
\text{H3 interaction required},\quad
\text{H3 atlas invalid},\quad
\text{H3 certificate incomplete}.
\end{gathered}
\]
\(\text{H3 stitch certified}\) is emitted only when every block above passes and the stable
assignment gap is positive. Equal-cost matchings, preserved nontrivial permutations of the visible
data, or coarse/fine gaps below margin make the fail-closed protocol emit
\(\text{H3 stitch ambiguous}\), meaning that it has certified no unique
stitch. This verdict is not by itself a theorem of mathematical nonuniqueness;
a symmetry obstruction requires a proved action that is fixed-point-free on
the admissible stitch set. Missing fields
force \(\text{H3 certificate incomplete}\). Failed atlas, interface, metric, or transport gates
force rejection or atlas invalidity. In particular, Euclidean distances on
the stored coordinate triples are not certificate distances; the certificate distance is
the hyperboloid geodesic \(d_H\) from Section~\ref{sec:h3-record-worldlines}.


\InputIfFileExists{tex_fragments/PRIOR_WORK_DERIVING_THE_PARTICLE_ZOO_FROM_OBSERVER_CONSISTENCY.tex}{}{}

\begin{thebibliography}{9}

\bibitem{ophrepo}
B.\ M\"uller et al.,
\emph{Observer-Patch Holography: formal proofs and executable receipts},
2026.
\url{https://github.com/FloatingPragma/observer-patch-holography/tree/r2054/Lean/Screen}

\bibitem{observableNormalForms}
B.\ M\"uller, J.\ Kim, D.\ Matscheko, and J.\ Hill,
\emph{Observation-Determined Normal Forms: Stability, Obstructions, and Refinement in Constraint
and Rewrite Systems},
2026.\\
Available at \url{https://github.com/FloatingPragma/observer-patch-holography/blob/r2054/extra/observable_normal_forms.pdf}.

\bibitem{ophsynthesis}
B.\ M\"uller, M.\ Poneder, A.\ Osika, K.\ Xue, B.\ Cassie, P.\ Nguyen, J.\ Kim, D.\ Matscheko, D.\ A.\ Friedman, J.\ Hill, W.\ T.\ Glynn, M.\ A.\ Visser, K.\ A.\ Anirudha, and B.\ de La Fourni\`ere,
\emph{Observers Who Read Their Universe into Existence: Finite Consensus, an Executable Observer, and Standard Physics}.
\url{https://github.com/FloatingPragma/observer-patch-holography/blob/r2054/flagship/from_observer_consensus_to_standard_physics.pdf}

\bibitem{ophobservers}
B.\ M\"uller, M.\ Poneder, A.\ Osika, K.\ Xue, B.\ Cassie, P.\ Nguyen, and D.\ A.\ Friedman,
\emph{Observers Are All You Need}.
GitHub PDF:
\url{https://github.com/FloatingPragma/observer-patch-holography/blob/r2054/paper/observers_are_all_you_need.pdf}

\bibitem{ophconsensus}
B.\ M\"uller, K.\ Xue, J.\ Kim, K.\ A.\ Anirudha, D.\ Matscheko, and J.\ Hill,
\emph{Reality as a Consensus Protocol: The Fixed-Point Computation That Implements Physics}.
GitHub PDF:
\url{https://github.com/FloatingPragma/observer-patch-holography/blob/r2054/paper/reality_as_consensus_protocol.pdf}

\bibitem{ophmicrophysics}
B.\ M\"uller, A.\ Osika, K.\ Xue, B.\ Cassie, B.\ de La Fourni\`ere, and D.\ A.\ Friedman,
\emph{Federated Echosahedral Screen Microphysics: Patch Hardware, Records, and Observer Synchronization in Observer-Patch Holography}.
GitHub PDF:
\url{https://github.com/FloatingPragma/observer-patch-holography/blob/r2054/paper/screen_microphysics_and_observer_synchronization.pdf}

\bibitem{ophfpe}
B.\ M\"uller,
\emph{OPH-FPE: finite simulator and receipt engine for Observer-Patch
Holography physics experiments},
2026.
Source:
\url{https://github.com/muellerberndt/oph-physics-sim}.
Canonical evidence archives:
\url{https://github.com/FloatingPragma/observer-patch-holography/tree/r2054/evidence}.
Exact local-domain archive and checker:
\url{https://github.com/FloatingPragma/observer-patch-holography/tree/r2054/evidence/local_domain}.

\bibitem{ophparticlesimdata}
FloatingPragma,
\emph{Particle-simulation receipt archive},
2026.
\url{https://github.com/FloatingPragma/observer-patch-holography/tree/r2054/evidence/particle_simulation_receipts}.

\bibitem{ophEinstein}
B.\ M\"uller, M.\ Poneder, A.\ Osika, K.\ Xue, P.\ Nguyen, M.\ A.\ Visser, D.\ Matscheko, and B.\ de La Fourni\`ere,
\emph{Recovering Observer Spacetime and Einstein Dynamics from Overlap Consistency}.
GitHub PDF:
\url{https://github.com/FloatingPragma/observer-patch-holography/blob/r2054/paper/recovering_observer_spacetime_and_einstein_dynamics_from_overlap_consistency.pdf}

\bibitem{ophGauge}
B.\ M\"uller, M.\ Poneder, A.\ Osika, K.\ Xue, P.\ Nguyen, D.\ Matscheko, and D.\ A.\ Friedman,
\emph{Standard Model Gauge Structure under Complete-Response and Observer-Transport Premises}.
GitHub PDF:
\url{https://github.com/FloatingPragma/observer-patch-holography/blob/r2054/paper/deriving_standard_model_gauge_structure_from_observer_overlap_consistency.pdf}

\bibitem{ophmodel}
B.\ M\"uller, M.\ Poneder, J.\ Kim, and D.\ A.\ Friedman,
\emph{An Executable Observer for Observer Patch Holography: Specification, Certified Components and Acceptance Criteria}.
GitHub PDF:
\url{https://github.com/FloatingPragma/observer-patch-holography/blob/r2054/paper/executable_observer_construction.pdf}

\bibitem{ophrepairlaw}
FloatingPragma,
\emph{Repair laws: what is fixed, what is free and how to tell},
2026.
\url{https://github.com/FloatingPragma/observer-patch-holography/blob/r2054/model/CANDIDATES.md}

\bibitem{codata2022}
National Institute of Standards and Technology,
\emph{2022 CODATA Recommended Values of the Fundamental Constants of Physics and Chemistry},
NIST SP 959, May 2024.
\url{https://physics.nist.gov/cuu/pdf/wallet_2022.pdf}

\bibitem{koide1983}
Y. Koide,
\emph{New view of quark and lepton mass hierarchy},
Physical Review D \textbf{28}, 252 (1983).

\bibitem{layton1990geometric}
E. Layton, Y. Huang, and S.-I. Chu,
\emph{Cyclic quantum evolution and Aharonov--Anandan geometric phases in
SU(2) spin-coherent states},
Physical Review A \textbf{41}, 42 (1990).
\url{https://doi.org/10.1103/PhysRevA.41.42}

\bibitem{rabei1999bargmann}
E. M. Rabei, Arvind, R. Simon, and N. Mukunda,
\emph{Bargmann invariants and geometric phases---a generalized connection},
Physical Review A \textbf{60}, 3397 (1999).
\url{https://arxiv.org/abs/quant-ph/9907080}

\bibitem{brannen2006}
C. A. Brannen,
\emph{The Lepton Masses},
preprint, 2006.
\url{http://brannenworks.com/MASSES2.pdf}

\bibitem{harremoesTopsoe2001}
P. Harremo\"es and F. Tops\o e,
\emph{Inequalities between entropy and index of coincidence derived from
information diagrams},
IEEE Transactions on Information Theory \textbf{47}, 2944--2960 (2001).
\url{https://doi.org/10.1109/18.959272}

\bibitem{antusch2025running}
S. Antusch, K. Hinze, and S. Saad,
\emph{Updated Running Quark and Lepton Parameters at Various Scales},
arXiv:2510.01312, 2025.
\url{https://arxiv.org/abs/2510.01312}

\bibitem{flag2024}
Y. Aoki et al. (Flavour Lattice Averaging Group),
\emph{FLAG Review 2024},
arXiv:2411.04268, 2024,
Tables 11--12.
\url{https://arxiv.org/abs/2411.04268}

\bibitem{pdg2024ckm}
S.~Navas et al. (Particle Data Group),
\emph{CKM Quark-Mixing Matrix},
Review of Particle Physics 2024, Sec.~12.2.2, revised May 31, 2024.
\url{https://pdg.lbl.gov/2024/reviews/rpp2024-rev-ckm-matrix.pdf}

\bibitem{pdg2025quarks}
S.~Navas et al. (Particle Data Group),
\emph{Review of Particle Physics}, Phys. Rev. D \textbf{110}, 030001 (2024),
2025 update, quark summary table.
\url{https://pdg.lbl.gov/2025/tables/rpp2025-sum-quarks.pdf}

\bibitem{buttazzo2013}
D.~Buttazzo et al.,
\emph{Investigating the near-criticality of the Higgs boson},
JHEP \textbf{12} (2013) 089, Appendix B.
\url{https://arxiv.org/abs/1307.3536}

\bibitem{pdg2026}
Particle Data Group,
\emph{Review of Particle Physics: 2026 particle listings},
2026.
\url{https://pdg.lbl.gov/2026/listings/particle_properties.html}

\bibitem{planck18}
Planck Collaboration,
\emph{Planck 2018 results. VI. Cosmological parameters},
Astronomy \& Astrophysics 641, A6, 2020.
\url{https://arxiv.org/abs/1807.06209}

\bibitem{desidr2nu}
DESI Collaboration,
\emph{Constraints on Neutrino Physics from DESI DR2 BAO and DR1 Full Shape},
arXiv:2503.14744, 2025.
\url{https://arxiv.org/abs/2503.14744}

\bibitem{nufit61}
I. Esteban, M. C. Gonzalez-Garcia, M. Maltoni, I. Martinez-Soler, J. P. Pinheiro, and T. Schwetz,
\emph{NuFit-6.0: Updated global analysis of three-flavor neutrino oscillations},
JHEP \textbf{12} (2024) 216, arXiv:2410.05380,
with the NuFIT 6.1 (2025) profile-table release at
\url{https://www.nu-fit.org/?q=node/309}.

\bibitem{kniehl2013}
B. A. Kniehl,
\emph{All-order renormalization of the propagator matrix for fermionic systems
with flavor mixing},
Physical Review D \textbf{89}, 096005 (2014), arXiv:1308.3140.

\bibitem{gambino1999}
P. Gambino and P. A. Grassi,
\emph{The Nielsen identities of the Standard Model and the definition of mass},
Physical Review D \textbf{62}, 076002 (2000), arXiv:hep-ph/9907254.

\bibitem{buchholz1986}
D. Buchholz,
\emph{Gauss' law and the infraparticle problem},
Physics Letters B \textbf{174}, 331--334 (1986).

\bibitem{duch2023}
P. Duch and W. Dybalski,
\emph{Infrared problem in quantum electrodynamics},
arXiv:2307.06114.

\bibitem{luscher1999}
M. L\"uscher,
\emph{Abelian chiral gauge theories on the lattice with exact gauge
invariance},
Nuclear Physics B \textbf{549}, 295--334 (1999), arXiv:hep-lat/9811032.

\bibitem{luscher2000}
M. L\"uscher,
\emph{Lattice regularization of chiral gauge theories to all orders of
perturbation theory},
Journal of High Energy Physics \textbf{06} (2000) 028, arXiv:hep-lat/0006014.

\bibitem{kadohKikukawa2008}
D. Kadoh and Y. Kikukawa,
\emph{A simple construction of fermion measure term in \(U(1)\) chiral
lattice gauge theories with exact gauge invariance},
Journal of High Energy Physics \textbf{02} (2008) 063, arXiv:0709.3658.

\bibitem{beluscaMaito2020}
H. B\'elusca-Ma\"ito, A. Ilakovac, M. Ma\dj or-Bo\v{z}inovi\'c, and
D. St\"ockinger,
\emph{Dimensional regularization and Breitenlohner--Maison/'t
Hooft--Veltman scheme for \(\gamma_5\) applied to chiral Yang--Mills theory},
Journal of High Energy Physics \textbf{08} (2020) 024, arXiv:2004.14398.

\bibitem{dudenasLoschner2020}
V. D\=ud\.enas and M. L\"oschner,
\emph{Vacuum expectation value renormalization in the Standard Model and
beyond},
Physical Review D \textbf{103}, 013003 (2021), arXiv:2010.15076.

\bibitem{dittmaierRzehak2022}
S. Dittmaier and H. Rzehak,
\emph{Electroweak renormalization based on gauge-invariant vacuum expectation
values of non-linear Higgs representations},
Journal of High Energy Physics \textbf{05} (2022) 125, arXiv:2203.07236.

\bibitem{grassiKniehlSirlin2001}
P. A. Grassi, B. A. Kniehl, and A. Sirlin,
\emph{Width and partial widths of unstable particles in the light of the
Nielsen identities},
Physical Review D \textbf{65}, 085001 (2002), arXiv:hep-ph/0109228.

\bibitem{MartinRobertsonSMDR}
S. P. Martin and D. G. Robertson,
\emph{Standard Model parameters in the tadpole-free pure
$\overline{\mathrm{MS}}$ scheme}, arXiv:1907.02500.

\bibitem{MartinTopPole}
S. P. Martin,
\emph{Top-quark pole mass in the tadpole-free MS-bar scheme},
arXiv:1604.01134.

\InputIfFileExists{tex_fragments/PRIOR_WORK_DERIVING_THE_PARTICLE_ZOO_FROM_OBSERVER_CONSISTENCY_BIBITEMS.tex}{}{}
\bibitem{BurgarthZeno2020}
D.~Burgarth, P.~Facchi, H.~Nakazato, S.~Pascazio and K.~Yuasa,
``Quantum Zeno Dynamics from General Quantum Operations,''
\emph{Quantum} \textbf{4}, 289 (2020),
\href{https://arxiv.org/abs/1809.09570}{arXiv:1809.09570}.

\bibitem{Susskind1977}
L.~Susskind, ``Lattice fermions,''
\href{https://doi.org/10.1103/PhysRevD.16.3031}{\emph{Phys. Rev. D}
\textbf{16}, 3031--3039 (1977)}.

\bibitem{MisumiStaggered2026}
T.~Misumi, T.~Onogi and T.~Yamaoka,
``Taste-splitting mass and edge modes in 3+1D staggered fermions,''
\href{https://arxiv.org/abs/2604.02078v1}{arXiv:2604.02078v1} (2026),
Sections~2.1--2.2.

\bibitem{CaputoExclusion2010}
P.~Caputo, T.~M.~Liggett and T.~Richthammer,
``Proof of Aldous' spectral gap conjecture,''
\emph{J. Amer. Math. Soc.} \textbf{23}, 831--851 (2010),
\href{https://arxiv.org/abs/0906.1238}{arXiv:0906.1238},
Theorem~1.1 and Section~4.1.1.

\end{thebibliography}

\section*{AI Assistance Disclosure}
This research project used research-grade commercial models, including
Anthropic's Fable and OpenAI's GPT-6 Astra, for research support, software
development, editing, and synthesis. The authors are responsible for the
paper's claims, methods, and final text.

\end{document}
